% ===========================================================================
%  Formal theory appendix.
%
%  Usage from main.tex (the caller decides the numbering):
%      \appendix
%      % ===========================================================================
%  Formal theory appendix.
%
%  Usage from main.tex (the caller decides the numbering):
%      \appendix
%      % ===========================================================================
%  Formal theory appendix.
%
%  Usage from main.tex (the caller decides the numbering):
%      \appendix
%      % ===========================================================================
%  Formal theory appendix.
%
%  Usage from main.tex (the caller decides the numbering):
%      \appendix
%      \input{theory_appendix}
%
%  The manuscript preamble loads amsmath, amssymb, and booktabs but not amsthm,
%  so results are set as \paragraph run-in headings and every proof ends with
%  \hfill$\square$. No preamble change is required, and this file can be built
%  on its own inside the same class.
%
%  Scope: this appendix proves only the statements used in the main text.
%  Values marked "P1 diagnostic" are registered outputs of
%  analysis/squad_qwen4b_p1/ and analysis/identity_continuous.json; they are
%  quoted, not re-estimated, and no gate is re-adjudicated here.
% ===========================================================================
\section{Formal theory of the depth decision}\label{app:theory}

This appendix states and proves the six results the manuscript relies on: the
first-order identity (Lemma~1), the shape condition that makes the identity
useful (Proposition~1), the comparative statics of the optimum (Corollary~1),
the boundary condition for a zero optimum (Corollary~2), the ceiling on the value
of per-question allocation (Proposition~2), and the cost of certifying a depth
(Proposition~3). Section~\ref{app:grid} separates the grid identity from the
continuous optimum, and Section~\ref{app:notclaimed} lists what the theory does
not license. Every result is conditional on the assumptions in
Section~\ref{app:assumptions}; none of them is a statement about the reader, the
retriever, or the hardware as such.

\subsection{Setting and assumptions}\label{app:assumptions}

Depth $k$ is a count of passages and ranges over the closed interval $[0,K]$ with
$K=32$. The measured grid is
\begin{equation}\label{app:eq:grid}
  \mathcal{K}=\{0,1,2,4,8,16,32\},
\end{equation}
so the grid is coarse and unevenly spaced, with gaps of $1$, $1$, $2$, $4$, $8$,
and $16$ passages. Mean energy per query is $J(k)$ in joules per query, accuracy
is $A(k)$ as a proportion, and energy per correct answer is the ratio estimand of
the main text,
\begin{equation}\label{app:eq:ratio}
  E(k)=\frac{J(k)}{A(k)}
  \qquad [\mathrm{J/correct}].
\end{equation}
The two elasticities and their gap are used throughout:
\begin{equation}\label{app:eq:elasticities}
  \varepsilon_{J}(k)=\frac{kJ'(k)}{J(k)},
  \qquad
  \varepsilon_{A}(k)=\frac{kA'(k)}{A(k)},
  \qquad
  \Delta(k)=\varepsilon_{J}(k)-\varepsilon_{A}(k),
\end{equation}
both dimensionless for $k>0$. The assumptions are as follows.

\begin{enumerate}
  \item[(A1)] \emph{Affine energy.} $J(k)=a+bk$ for $k\in[0,K]$, with $a>0$ in
  J/query and $b>0$ in J per passage per query. Equivalently
  $\varepsilon_{J}(k)=bk/(a+bk)\in(0,1)$ and $\varepsilon_{J}$ is strictly
  increasing in $k$.

  \item[(A2)] \emph{Positive accuracy.} $A(k)\in(0,1]$ for $k\in[0,K]$, and $A$ is
  twice differentiable wherever a derivative appears. Curvature signs are stated
  per result and are never assumed globally without being named.

  \item[(A3)] \emph{Decision object.} The policy quantity is the constrained
  minimiser over the grid, $k^{*}_{m}=\arg\min_{k\in\mathcal{F}_{m}}E(k)$ with
  $\mathcal{F}_{m}=\{k\in\mathcal{K}:A(k)\ge A(K)-m\}$, which is the object of the
  main text. Statements about a continuous $k^{*}$ are statements about the
  extension in (A4), not about unmeasured depths.

  \item[(A4)] \emph{Continuum extension.} Off the grid, $A$ denotes a twice
  differentiable interpolant of the seven grid means (the registered analysis
  uses a monotone cubic). The interpolant is a modelling device: it makes
  $E'(k)$ well defined, and every off-grid number inherits its smoothness
  assumption.

  \item[(A5)] \emph{Sampling.} Questions are independent draws from a fixed
  distribution, each question is answered at every depth in $\mathcal{K}$, and
  correctness is binary, $C_{qk}\in\{0,1\}$ with $\mathbb{E}[C_{qk}]=A(k)$. This
  assumption is used only in Proposition~\ref{app:prop:cert}(b).

  \item[(A6)] \emph{Calibration.} Calibrated energy is a positive depth-common
  multiple of the measured energy, $J_{\mathrm{cal}}(k)=fJ(k)$ with $f>0$. The
  calibration rule of the main text is all-or-nothing across depths, so this is
  the regime in which absolute energy is reported.
\end{enumerate}

Assumption (A6) has an immediate consequence that justifies using calibrated
energy in the theory: $E_{\mathrm{cal}}=fE$, so the minimiser, every ratio of two
depths, and every elasticity $\varepsilon_{J}$ are unchanged, because $f$ cancels
in $kJ'_{\mathrm{cal}}/J_{\mathrm{cal}}$. A depth-dependent calibration error
would not cancel, which is why the measurement protocol gates on a per-depth
factor and applies it only if all depths pass.

\begin{table}[pos=htbp]
  \centering
  \caption{Summary of the appendix. Each result is conditional on the
  hypotheses named in its statement; none of them is a claim about a particular
  reader, corpus, or GPU.}
  \label{tab:theorymap}
  \footnotesize
  \setlength{\tabcolsep}{4pt}
  \begin{tabular}{@{}>{\raggedright\arraybackslash}p{1.9cm}>{\raggedright\arraybackslash}p{7.4cm}>{\raggedright\arraybackslash}p{5.0cm}@{}}
    \toprule
    \textbf{Result} & \textbf{What is proved} & \textbf{What it licenses} \\
    \midrule
    Lemma~1 & $E'(k)=0$ iff the two elasticities agree, for $k>0$ &
    A local diagnostic; necessary, not sufficient \\
    Proposition~1 & Concave $A$ makes the stationarity function monotone, so $E$
    is unimodal & Existence and uniqueness of an interior optimum, and a test of
    the hypothesis \\
    Corollary~1 & $\partial k^{*}/\partial a>0$ and $\partial k^{*}/\partial b<0$
    at a slack-gate interior optimum & Directional, qualitative predictions only \\
    Corollary~2 & A closed-form $b/a$ threshold for $k^{*}=0$ under coverage &
    Boundary optima are possible; a check on the coverage family only \\
    Proposition~2 & Any allocation at matched accuracy satisfies
    $E_{\pi}/E_{s}\le 1+bK/a$ & A ceiling on the value of adaptivity, with a
    sharpness example \\
    Proposition~3 & Break-even volume and sample complexity of certification &
    A cost, not a guarantee; the accuracy side dominates \\
    \bottomrule
  \end{tabular}
\end{table}

\subsection{Lemma 1: the first-order identity}\label{app:lemma1}

\paragraph{Lemma 1.} Under (A1) and (A2), for every $k\in(0,K]$,
\begin{equation}\label{app:eq:identity}
  E'(k)=\frac{bA(k)-J(k)A'(k)}{A(k)^{2}}
  =\frac{J(k)}{kA(k)}\,\Delta(k),
\end{equation}
and therefore
\begin{equation}\label{app:eq:identity-iff}
  E'(k)=0
  \iff
  bA(k)=J(k)A'(k)
  \iff
  \frac{kJ'(k)}{J(k)}=\frac{kA'(k)}{A(k)}
  \iff
  \Delta(k)=0 .
\end{equation}
At $k=0$ only the unmultiplied form is meaningful,
$E'(0)=[bA(0)-aA'(0)]/A(0)^{2}$, because the elasticity form is identically
$0=0$ there.

\paragraph{Proof.} By the quotient rule and (A1),
\begin{equation}\label{app:eq:proof-quotient}
  E'(k)=\frac{J'(k)A(k)-J(k)A'(k)}{A(k)^{2}}
  =\frac{bA(k)-J(k)A'(k)}{A(k)^{2}},
\end{equation}
which is the first equality in \eqref{app:eq:identity}; $A(k)>0$ by (A2), so the
expression is defined. For the second equality, write
$E=J/A$ and differentiate logarithmically, $E'/E=J'/J-A'/A$. Multiplying by $k$
gives $kE'(k)/E(k)=\Delta(k)$ and hence
\begin{equation}\label{app:eq:proof-second}
  E'(k)=\frac{E(k)}{k}\Delta(k)=\frac{J(k)}{kA(k)}\Delta(k),
\end{equation}
which agrees with \eqref{app:eq:proof-quotient} after multiplying both sides by
$kA(k)^{2}$. Since $J(k)>0$, $A(k)>0$, and $k>0$, the prefactor $J/(kA)$ in
\eqref{app:eq:proof-second} is strictly positive, so $E'(k)=0$ holds if and only
if $\Delta(k)=0$; the middle equivalence in \eqref{app:eq:identity-iff} is the
same statement multiplied by $k>0$ and divided by the positive quantities $J$
and $A$. \hfill$\square$

\paragraph{Necessity, not sufficiency.} An interior local extremum of a
differentiable $E$ must satisfy \eqref{app:eq:identity-iff}, so the identity is a
necessary condition for an interior optimum and it identifies the candidates.
It is not sufficient: a point where the elasticities agree may be a local
maximum, and the identity says nothing about a boundary optimum. The smallest
counterexample uses a convex accuracy curve, $a=b=1$ and
$A(k)=0.1+0.05k^{2}$. Then $\Delta(k)=0$ at $k=\sqrt{3}-1\approx0.732$, where
$E\approx13.66$ J/correct, while $E(0)=E(2)=10$ J/correct; the stationary point
is a local maximum, and the hypothesis $A''\le0$ of Proposition~\ref{app:prop:unimodal}
fails. A zero gap is therefore evidence about the local slope of $E$ at a depth,
never by itself evidence that the depth is optimal.

\paragraph{Reading the sign.} Because the prefactor in \eqref{app:eq:proof-second}
is positive, $\operatorname{sign}E'(k)=\operatorname{sign}\Delta(k)$ for $k>0$. A
positive gap means the marginal joule buys less than the marginal correct answer,
so a small increase in depth raises energy per correct answer; a negative gap
means the opposite. This is the interpretation used for the reported gap, and it
is invariant to the depth-common factor $f$ of (A6), because both elasticities
are logarithmic derivatives.

\subsection{Proposition 1: unimodality under concave accuracy}\label{app:prop:unimodal}

\paragraph{Proposition 1.} Assume (A1), (A2), and $A''(k)\le0$ on $(0,K)$; write
$J(k)=a+bk$ and define the stationarity function
\begin{equation}\label{app:eq:N}
  N(k)=bA(k)-J(k)A'(k)=A(k)^{2}E'(k)
  \qquad [\mathrm{J\,query^{-1}\,passage^{-1}}].
\end{equation}
Then:
\begin{enumerate}
  \item[(i)] $N'(k)=-(a+bk)A''(k)\ge0$, so $N$ is nondecreasing on $[0,K]$, and
  strictly increasing if $A''<0$ on $(0,K)$.
  \item[(ii)] $E'=N/A^{2}$ changes sign at most once on $(0,K)$, and any sign
  change is from negative to positive. Hence $E$ is nonincreasing then
  nondecreasing: it is unimodal (quasiconvex) and has no interior local maximum.
  If $A''<0$, $E$ has at most one interior stationary point, and when it exists it
  is the unique global minimiser of $E$ on $[0,K]$.
  \item[(iii)] If $A''<0$, an interior minimiser exists if and only if
  \begin{equation}\label{app:eq:interior}
    bA(0)<aA'(0)
    \quad\text{and}\quad
    bA(K)>J(K)A'(K);
  \end{equation}
  otherwise the minimiser is $k=0$ (if the first inequality fails) or $k=K$ (if
  the second fails).
\end{enumerate}

\paragraph{Proof.} (i) Differentiate \eqref{app:eq:N} using (A1):
\begin{equation}\label{app:eq:Nprime}
  N'(k)=bA'(k)-\left[bA'(k)+(a+bk)A''(k)\right]=-(a+bk)A''(k),
\end{equation}
and $a+bk>0$ by (A1), so $N'\ge0$; if $A''<0$ then $N'>0$.

(ii) By \eqref{app:eq:N}, $E'=N/A^{2}$ and $A^{2}>0$, so $E'$ has the sign of
$N$. A nondecreasing function crosses zero at most once, and when it does the
crossing is from negative to positive; hence there is a point $\bar k$ (possibly
$0$ or $K$, possibly an interval when $N$ is flat at zero) with $E'\le0$ for
$k<\bar k$ and $E'\ge0$ for $k>\bar k$. That is unimodality, and no interior local
maximum is possible. If $A''<0$ then $N$ is strictly increasing, so its zero set
is empty or a single point $k^{*}$. At such a point,
\begin{equation}\label{app:eq:Esecond}
  E''(k^{*})=\frac{N'(k^{*})A(k^{*})^{2}-2A(k^{*})A'(k^{*})N(k^{*})}{A(k^{*})^{4}}
  =\frac{N'(k^{*})}{A(k^{*})^{2}}>0,
\end{equation}
because $N(k^{*})=0$; the point is a strict local minimum, and by the sign pattern
it is the unique global minimiser on $[0,K]$.

(iii) If $A''<0$ then $N$ is continuous and strictly increasing. An interior zero
exists if and only if $N(0)<0<N(K)$; substituting \eqref{app:eq:N} gives
$bA(0)-aA'(0)<0$ and $bA(K)-(a+bK)A'(K)>0$, which is \eqref{app:eq:interior}.
If $N(0)\ge0$ then $N>0$ on $(0,K]$ and $E$ is nondecreasing, so the minimiser is
$k=0$; if $N(K)\le0$ then $N<0$ on $[0,K)$ and $E$ is nonincreasing, so the
minimiser is $k=K$. \hfill$\square$

\paragraph{Two remarks.} First, the proof never uses $A'\ge0$: monotonicity of
$N$ depends only on the sign of $A''$. The proposition therefore also covers a
concave decline after an accuracy peak, which is the qualitative shape of the two
non-monotone P4 arms (a peak followed by a decline), although concavity of the
decline is not tested there. Second, the hypothesis is testable on the grid: the discrete
analogue of $A''\le0$ is that the secant slopes between adjacent grid depths are
non-increasing, which is what the companion implementation checks. The hypothesis
is not a description of every accuracy curve; where it fails, Proposition~1 gives
no global statement and the identity remains local.

\paragraph{Existence in elasticity form.} Conditions \eqref{app:eq:interior} can
be read as bounds on the elasticities. The first is
$b/a<A'(0)/A(0)$, i.e.\ the relative accuracy gain per passage at zero must exceed
the relative energy cost per passage. The second is
$\varepsilon_{A}(K)<\varepsilon_{J}(K)=bK/(a+bK)<1$, i.e.\ the accuracy
elasticity at the top of the grid must be below the energy elasticity. Both are
properties of the fitted curves, not of the decision rule, and both are evaluated
without the accuracy gate of (A3).

\subsection{Corollary 1: comparative statics of the optimum}\label{app:cor:statics}

\paragraph{Corollary 1.} Assume (A1), (A2), $A'(k^{*})>0$ and $A''(k^{*})<0$ at an
interior stationary point $k^{*}=k^{*}(a,b)$ of $E$. Then $k^{*}$ is continuously
differentiable in $(a,b)$ in a neighbourhood, and
\begin{equation}\label{app:eq:statics}
  \frac{\partial k^{*}}{\partial a}
  =\frac{-A'(k^{*})}{(a+bk^{*})A''(k^{*})}>0,
  \qquad
  \frac{\partial k^{*}}{\partial b}
  =\frac{A(k^{*})-k^{*}A'(k^{*})}{(a+bk^{*})A''(k^{*})}
  =-\frac{aA(k^{*})}{(a+bk^{*})^{2}\left(-A''(k^{*})\right)}<0 .
\end{equation}

\paragraph{Proof.} By Lemma~1 the stationary point solves
$F(k,a,b)=bA(k)-(a+bk)A'(k)=0$, and $F$ is continuously differentiable under
(A2). Its partials are
\begin{equation}\label{app:eq:Fpartials}
  F_{k}=-(a+bk)A''(k),
  \qquad
  F_{a}=-A'(k),
  \qquad
  F_{b}=A(k)-kA'(k).
\end{equation}
At $k^{*}$, $F_{k}=-(a+bk^{*})A''(k^{*})>0$ because $A''(k^{*})<0$; this is also
the second-order condition \eqref{app:eq:Esecond}. The implicit function theorem
therefore applies, and
\begin{equation}\label{app:eq:IFT}
  \frac{\partial k^{*}}{\partial a}=-\frac{F_{a}}{F_{k}},
  \qquad
  \frac{\partial k^{*}}{\partial b}=-\frac{F_{b}}{F_{k}} .
\end{equation}
Substituting \eqref{app:eq:Fpartials} gives the first equality in each line of
\eqref{app:eq:statics}. The first sign is immediate since $A'(k^{*})>0$ and the
denominator is negative. For the second, apply the identity
\eqref{app:eq:identity-iff} at $k^{*}$:
$k^{*}A'(k^{*})/A(k^{*})=\varepsilon_{J}(k^{*})=bk^{*}/(a+bk^{*})<1$, hence
\begin{equation}\label{app:eq:Fbpositive}
  A(k^{*})-k^{*}A'(k^{*})
  =A(k^{*})\left(1-\varepsilon_{J}(k^{*})\right)
  =\frac{aA(k^{*})}{a+bk^{*}}>0,
\end{equation}
and the denominator of the second line is negative, so
$\partial k^{*}/\partial b<0$. \hfill$\square$

\paragraph{Reading.} A larger fixed cost per query makes depth relatively cheap
and pushes the unconstrained optimum deeper; a larger per-passage cost makes
depth relatively expensive and pushes it shallower. The corollary is local and
applies only when the accuracy gate is slack: when the gate binds, the policy
depth is the smallest feasible grid depth and is locally insensitive to $(a,b)$,
so the derivative is zero in the one-sided sense at that corner. It is also
silent about the magnitude, because $A'(k^{*})$ and $A''(k^{*})$ are estimated
from the same finite sample as $A(k^{*})$; the sign is a prediction, the size is
not.

\subsection{Corollary 2: when the optimum is the boundary}\label{app:cor:boundary}

\paragraph{Corollary 2.} Assume (A1) and the stylized coverage family
\begin{equation}\label{app:eq:coverage}
  A(k)=1-(1-a_{0})(1-p)^{k},
  \qquad
  a_{0}=A(0)\in(0,1),
  \qquad
  p\in(0,1),
\end{equation}
so that $A(k)\in(a_{0},1)$ and (A2) holds on $[0,\infty)$. Let
$L=-\ln(1-p)>0$. Then $A'(k)=L\left(1-A(k)\right)>0$ and
$A''(k)=-L^{2}\left(1-A(k)\right)<0$, Proposition~\ref{app:prop:unimodal}
applies, and
\begin{equation}\label{app:eq:boundary}
  k^{*}=0
  \iff
  \frac{b}{a}\ge\frac{(1-a_{0})\left(-\ln(1-p)\right)}{a_{0}}
  =\frac{1-a_{0}}{a_{0}}\,L .
\end{equation}
When the inequality is reversed, the unique interior optimum solves
\begin{equation}\label{app:eq:coverage-optimum}
  (1-a_{0})(1-p)^{k^{*}}=\frac{b}{aL+b+bLk^{*}} .
\end{equation}

\paragraph{Proof.} Differentiating \eqref{app:eq:coverage} gives
$A'(k)=-(1-a_{0})(1-p)^{k}\ln(1-p)=L\left(1-A(k)\right)$ and
$A''(k)=-L^{2}\left(1-A(k)\right)<0$, so $A$ is strictly increasing and strictly
concave and (A2) holds. Evaluate \eqref{app:eq:N} at zero:
\begin{equation}\label{app:eq:N0}
  N(0)=bA(0)-aA'(0)=ba_{0}-a(1-a_{0})L
  =a_{0}\left[b-a\frac{1-a_{0}}{a_{0}}L\right].
\end{equation}
By Proposition~\ref{app:prop:unimodal}(i), $N$ is strictly increasing. As
$k\to\infty$, $A(k)\to1$ and $(a+bk)(1-A(k))\to0$ because the exponential decay
dominates the linear factor, so
\begin{equation}\label{app:eq:Nlimit}
  N(k)=bA(k)-(a+bk)L\left(1-A(k)\right)\longrightarrow b>0 .
\end{equation}
If $N(0)\ge0$, then $N>0$ on $(0,\infty)$ and $E$ is strictly increasing, so
$k^{*}=0$; by \eqref{app:eq:N0} this is exactly the inequality in
\eqref{app:eq:boundary}. If $N(0)<0$, then \eqref{app:eq:Nlimit} and strict
monotonicity give a unique zero $k^{*}\in(0,\infty)$, which is the unique
minimiser by Proposition~\ref{app:prop:unimodal}(ii). Setting $N(k^{*})=0$ and
writing $s=1-A(k^{*})=(1-a_{0})(1-p)^{k^{*}}$ gives
$b(1-s)=(a+bk^{*})Ls$, i.e.\ $b=s\left(aL+b+bLk^{*}\right)$, which is
\eqref{app:eq:coverage-optimum}. \hfill$\square$

\paragraph{Scope.} The corollary is a statement about one stylized family, not
about the observed curves. The P4 validation rejects the family as a general
description of the four arms ($R^{2}$ between $0.097$ and $0.879$), so its role
is to show that a boundary optimum is a possible outcome of diminishing returns,
and to give the threshold in closed form. It does not certify any observed
$k^{*}$, and it is not used to choose a depth.

\subsection{Proposition 2: the price of static depth}\label{app:prop:price}

\paragraph{Proposition 2.} Assume (A1). Let $\pi$ be any per-question depth
assignment with values in $[0,K]$, with population energy
$J_{\pi}=\mathbb{E}_{q}[J(\pi(q))]$ and accuracy
$A_{\pi}=\mathbb{E}_{q}[A(\pi(q))]$, and let $k_{s}$ be a static depth with
accuracy $A_{s}=A(k_{s})$. If $A_{\pi}\ge A_{s}$, then
\begin{equation}\label{app:eq:ceiling}
  \frac{E_{\pi}}{E_{s}}
  =\frac{J_{\pi}}{J(k_{s})}\cdot\frac{A_{s}}{A_{\pi}}
  \le\frac{a+bK}{a}
  =\frac{1}{1-\varepsilon_{J}(K)} ,
\end{equation}
and the largest saving any such assignment can deliver is the variable energy
share at the reference depth,
\begin{equation}\label{app:eq:ceiling-saving}
  1-\frac{E_{\pi}}{E_{s}}\le\frac{bK}{a+bK}=\varepsilon_{J}(K)<1 .
\end{equation}
The bound is sharp. For $\varepsilon\in(0,1)$ let a fraction $1-\varepsilon$ of
questions be easy, answered correctly at every depth including $k=0$, and a
fraction $\varepsilon$ be hard, answered correctly only at depths $k\ge K$. Then
the best static depth at accuracy $1$ is $k_{s}=K$, the hindsight oracle
$\pi_{\mathrm{or}}$ assigns easy questions to $0$ and hard questions to $K$, and
\begin{equation}\label{app:eq:rho}
  \rho(\varepsilon)
  =\frac{E_{\mathrm{static}}}{E_{\mathrm{oracle}}}
  =\frac{a+bK}{a+\varepsilon bK}
  \;\longrightarrow\;
  1+\frac{bK}{a}
  =\frac{1}{1-\varepsilon_{J}(K)}
  \qquad(\varepsilon\to0).
\end{equation}
The ratio is written here as static over oracle, so that the quantity is the
\emph{price} of fixing one depth; written in the opposite order (oracle over
static) the same limit is its reciprocal $a/(a+bK)<1$.

\paragraph{Proof.} By (A1), $J(\pi(q))\in[a,a+bK]$ for every question, so
$a\le J_{\pi}\le a+bK$ and $J(k_{s})\ge a$; hence
$J_{\pi}/J(k_{s})\le(a+bK)/a$. Since $A_{\pi}\ge A_{s}>0$, the factor
$A_{s}/A_{\pi}$ is at most one, which proves \eqref{app:eq:ceiling}. For
\eqref{app:eq:ceiling-saving}, subtract the reciprocal:
$1-a/(a+bK)=bK/(a+bK)=\varepsilon_{J}(K)$ by (A1).

For sharpness, let $q$ be easy with probability $1-\varepsilon$ and hard with
probability $\varepsilon$. Every static depth $k<K$ answers only the easy
questions, so its accuracy is $1-\varepsilon<1$; the only static depths with
accuracy $1$ are $k\ge K$, and among them $J$ is minimised at $k=K$. Hence
$E_{\mathrm{static}}=a+bK$. The oracle answers every question correctly, so its
accuracy is $1$ and its energy is
$(1-\varepsilon)a+\varepsilon(a+bK)=a+\varepsilon bK$, giving
$E_{\mathrm{oracle}}=a+\varepsilon bK$ and the first equality in
\eqref{app:eq:rho}. The limit is immediate, and the last equality repeats
$bK/(a+bK)=\varepsilon_{J}(K)$. \hfill$\square$

\paragraph{Scope.} Proposition~\ref{app:prop:price} is an upper bound with a
sharpness example, not a lower bound on the value of adaptivity and not a claim
about the observed workload. Two consequences matter for the paper. First, the
ceiling is a property of the affine energy model alone: it caps \emph{any}
allocation, learned or oracle, at matched accuracy, and it is the formal reason
the oracle gain is reported as a fraction of the static energy rather than as an
unbounded promise. Second, the example shows that the ceiling is approached when
a small hard stratum forces the whole workload to the deepest setting; when the
stratum is large ($\varepsilon\to1$), the ratio tends to one and adaptivity is
worth nothing. The empirical matched-oracle gain reported for P3 is far below the
P1 ceiling of $\varepsilon_{J}(32)=0.7517$, as it must be, and the pilot oracle
figure of $78.8\%$ is below the pilot ceiling
$bK/(a+bK)=0.8227$ computed from $a=15.74$ and $b=2.283$; both are consistency
checks on registered quantities, not new evidence.

\subsection{Proposition 3: the cost of certification}\label{app:prop:cert}

Certification has two separate costs: the energy of the experiment (part a) and
the number of questions needed for the statistical statement (part b). They are
stated separately because the first is a deterministic break-even and the second
is a sampling requirement, and only their product is the cost of a certified
depth.

\paragraph{Proposition 3(a): economic break-even.} Assume (A1), a deployment of
$N$ queries drawn from the same distribution as the experiment, one depth used
for all $N$ queries, no discounting, and a one-off certification cost
$C_{\mathrm{cert}}>0$ in joules. Let $k_{0}=K$ be the reference depth and
$k^{*}$ the candidate. Define the per-query break-even and the
per-correct-answer break-even by
\begin{equation}\label{app:eq:nstar}
  N^{*}_{J}=\frac{C_{\mathrm{cert}}}{J(k_{0})-J(k^{*})},
  \qquad
  N^{*}_{A}=\frac{C_{\mathrm{cert}}}{A(k^{*})\left[E(k_{0})-E(k^{*})\right]},
\end{equation}
both in queries and both defined only when the denominator is positive. Then
\begin{equation}\label{app:eq:nstar-identity}
  A(k^{*})\left[E(k_{0})-E(k^{*})\right]
  =\frac{A(k^{*})}{A(k_{0})}J(k_{0})-J(k^{*}),
\end{equation}
so $N^{*}_{A}=N^{*}_{J}$ if and only if $A(k^{*})=A(k_{0})$, and otherwise
\begin{equation}\label{app:eq:beta}
  \beta:=\frac{N^{*}_{A}}{N^{*}_{J}}
  =\frac{J(k_{0})-J(k^{*})}{\dfrac{A(k^{*})}{A(k_{0})}J(k_{0})-J(k^{*})},
\end{equation}
with $\beta\ge1$ if and only if $A(k^{*})\le A(k_{0})$.

\paragraph{Proof.} Equation \eqref{app:eq:nstar-identity} is
$A(k^{*})E(k_{0})-A(k^{*})E(k^{*})=A(k^{*})J(k_{0})/A(k_{0})-J(k^{*})$, using
$A(k^{*})E(k^{*})=J(k^{*})$. Dividing the two lines of \eqref{app:eq:nstar}
gives \eqref{app:eq:beta}. For the sign, compare denominators:
\begin{equation}\label{app:eq:beta-sign}
  \left[J(k_{0})-J(k^{*})\right]
  -\left[\frac{A(k^{*})}{A(k_{0})}J(k_{0})-J(k^{*})\right]
  =J(k_{0})\left(1-\frac{A(k^{*})}{A(k_{0})}\right),
\end{equation}
which is non-negative exactly when $A(k^{*})\le A(k_{0})$. Both denominators are
positive by hypothesis, so the inequality passes to $\beta$.
\hfill$\square$

\paragraph{Interpretation and the approximation.} $N^{*}_{J}$ is the number of
deployment queries at which the energy saved by switching depth equals the
certificate: after $N$ queries at $k^{*}$ instead of $k_{0}$ the saving is
$N\left[J(k_{0})-J(k^{*})\right]$, and it repays $C_{\mathrm{cert}}$ at
$N=N^{*}_{J}$. $N^{*}_{A}$ answers a different question: it values the candidate's
correct-answer output at the per-correct-answer saving $E(k_{0})-E(k^{*})$, so $N$
candidate-depth queries produce $N A(k^{*})$ correct answers and a cumulative
saving $N A(k^{*})\left[E(k_{0})-E(k^{*})\right]$, which equals
$C_{\mathrm{cert}}$ at $N=N^{*}_{A}$. The two coincide only at
equal accuracy. When $A(k^{*})<A(k_{0})$, the paper's per-correct-answer formula
is therefore an approximation to the deployment-queries reading, and it errs on
the conservative side ($\beta>1$): it reports a longer break-even because it also
requires the candidate to reproduce the correct-answer output that its lower
accuracy would otherwise lose. The direction is known but the formula is not
exact, which is why both numbers are reported rather than one.

\paragraph{P1 diagnostic.} With the registered P1 values
$C_{\mathrm{cert}}=311{,}992.064$ J, $J(4)=36.2537$ J/query,
$J(32)=94.9450$ J/query, $A(4)=0.450$, and $A(32)=0.493$, one obtains
$N^{*}_{J}=5{,}315.8$ queries, $N^{*}_{A}=6{,}189.1$ queries, and
$\beta=1.1643$. The certification cost is itself a sum over the grid,
$C_{\mathrm{cert}}=n\sum_{k\in\mathcal{K}}J(k)=1{,}000\times311.992$ J, so the
sweep energy and the reported break-even use the same accounting.

\paragraph{Proposition 3(b): statistical sample complexity.} Assume (A5), let
$m$ be the accuracy margin (written $\delta$ in the main text, and renamed here
to avoid a clash with the error probability), and let
$\Delta=A(k_{0})-A(k^{*})>0$. For the paired differences
$D_{q}=C_{qk_{0}}-C_{qk^{*}}\in\{-1,0,1\}$ with
$\mathbb{E}[D_{q}]=\Delta$, put
$t_{n}=\sqrt{2\ln(1/\alpha)/n}$ and declare non-inferiority of $k^{*}$ when
\begin{equation}\label{app:eq:declare}
  \bar D+t_{n}<m .
\end{equation}
Then:
\begin{enumerate}
  \item[(i)] \eqref{app:eq:declare} is a level-$\alpha$ test of
  $H_{0}:\Delta\ge m$.
  \item[(ii)] At a true difference $\Delta=m_{1}<m$ the probability of failing to
  declare non-inferiority is at most
  \begin{equation}\label{app:eq:hoeffding}
    \Pr\left(\text{fail to declare}\mid \Delta=m_{1}\right)
    \le\exp\left(-\frac{n\left(\left(m-m_{1}\right)-t_{n}\right)^{2}}{2}\right),
  \end{equation}
  so $n\ge8\max\left(\ln(1/\alpha),\ln(1/\delta_{\mathrm{err}})\right)/(m-m_{1})^{2}$
  paired questions suffice for a failure probability at most
  $\delta_{\mathrm{err}}$.
  \item[(iii)] Selecting the more accurate of two depths at a true gap $\Delta$
  needs $n\ge2\ln(1/\delta_{\mathrm{err}})/\Delta^{2}$, and selecting among $d$
  depths whose smallest gap to the optimum is $\Delta_{\min}$ needs
  $n\ge2\ln\left((d-1)/\delta_{\mathrm{err}}\right)/\Delta_{\min}^{2}$ by a union
  bound; in both cases the correct depth is selected with probability at least
  $1-\delta_{\mathrm{err}}$.
\end{enumerate}

\paragraph{Proof sketch.} $D_{q}$ is bounded in $[-1,1]$ and has range $2$, so
Hoeffding's inequality gives, for every $u>0$,
\begin{equation}\label{app:eq:two-tails}
  \Pr\left(\bar D-\Delta\le-u\right)\le\exp\left(-\frac{nu^{2}}{2}\right),
  \qquad
  \Pr\left(\bar D-\Delta\ge u\right)\le\exp\left(-\frac{nu^{2}}{2}\right).
\end{equation}
(i) Under $H_{0}$, $\Delta\ge m$, so declaring non-inferiority implies
$\bar D-\Delta<-t_{n}$, which has probability at most
$\exp(-nt_{n}^{2}/2)=\alpha$.

(ii) Failure to declare means $\bar D+t_{n}\ge m$, hence
$\bar D-\Delta\ge(m-m_{1})-t_{n}$, and the upper tail of
\eqref{app:eq:two-tails} applies. If $t_{n}\le(m-m_{1})/2$, which holds once
$n\ge8\ln(1/\alpha)/(m-m_{1})^{2}$, the exponent is at most
$-n(m-m_{1})^{2}/8$, and the stated $n$ makes it at most
$\ln(1/\delta_{\mathrm{err}})$.

(iii) The selection error is
$\Pr(\bar D\le0)=\Pr(\bar D-\Delta\le-\Delta)\le\exp(-n\Delta^{2}/2)$, and the
union bound over the $d-1$ alternatives gives the second statement.
\hfill$\square$

\paragraph{Scope and consequences.} The bound is distribution-free and makes no
normal approximation, but it is deliberately conservative: it ignores the
pairing, i.e.\ the discordance rate, which is what a McNemar calculation uses.
The tighter normal-approximation requirement scales like
$\psi/(m-m_{1})^{2}$ at discordance $\psi$, so at a fixed difference the observed
P1 discordance of $13.1\%$ requires roughly $1.6$ times the pairs that the
planned $8\%$ implied, and the registered frontier puts the 5 pp requirement at
$n=7{,}138$ pairs. Three consequences are used in the main text. First, the
sample complexity diverges as $m_{1}\uparrow m$: a margin that the data show is
several percentage points wide cannot be resolved by a design powered at a
near-zero difference, which is the formal content of the P1 precision failure.
Second, the energy side does not drive the requirement: per-query energy has a P1
repeat CV of $2.81\%$ at $k=8$, so the relative standard error of a mean energy
at $n=1{,}000$ is about $2.81\%/\sqrt{1{,}000}<0.1\%$, while the accuracy
contrast at 5 pp remains unresolved at the same $n$. The accuracy side is the
binding one. Third, the certification cost is linear in the statistical
requirement, because a sweep of $n$ questions costs
$C_{\mathrm{cert}}(n)=n\,c_{\mathrm{sweep}}$ with
$c_{\mathrm{sweep}}=\sum_{k\in\mathcal{K}}J(k)$ in J/query, so
\begin{equation}\label{app:eq:nstar-n}
  N^{*}_{J}(n)=\frac{n\,c_{\mathrm{sweep}}}{J(k_{0})-J(k^{*})},
\end{equation}
which grows linearly in $n$. The statistical requirement and the economic
break-even therefore cannot be chosen independently, and a depth that needs a
large experiment needs a proportionally larger deployment to justify it.

\subsection{Grid resolution and the H3 diagnostic}\label{app:grid}

The identity of Lemma~1 is a statement about a derivative. On the grid two
approximations intervene, and separating them is what makes the H3 result
readable.

\emph{First, the derivative is a secant.} The reported accuracy elasticity at a
grid depth is a finite difference over the neighbouring depths, so the identity
at a grid point is the statement that the secant elasticity equals the affine
elasticity. At $k=4$ the central difference uses $k=2$ and $k=8$, a step of three
passages; at $k=8$ it uses $k=4$ and $k=16$, a step of six; at the edges the
difference is one-sided, with step one at $k=1$ and step sixteen at $k=32$. A
central difference approximates $A'(k)$ to second order in the step and a
one-sided difference to first order, so a grid gap is a discretisation of
$\Delta(k)$, not a different quantity.

\emph{Second, the decision object is constrained.} The policy depth of (A3) is the
energy minimiser inside the feasible set $\mathcal{F}_{m}$, whereas $k^{*}$ in
Lemma~1 is the unconstrained stationary point of $E$. The two objects coincide
only when the accuracy gate is slack. In P1 the gate binds: the 5 pp rule selects
$k=4$ because it is the smallest feasible depth, not because $E$ is stationary
there. Testing the identity at the gate-selected depth is therefore a test of
local stationarity at a point that the theory does not expect to be stationary,
and the main text reports it as such.

\emph{The P1 numbers.} With $a=24.0759$ J/query and
$b=2.2772$ J per passage per query, the affine elasticity is
$\varepsilon_{J}(4)=0.2745$, the secant elasticity is
$\varepsilon_{A}(4)=0.0741$, and the gap is $\Delta(4)=+0.2004$ with a 95\%
interval of $[0.1677,0.2320]$ that excludes zero. The positive sign says that
$E$ is increasing at $k=4$: the last passages bought at that depth cost more
energy per correct answer than they return. The monotone-cubic interpolant of the
same seven means has its unconstrained stationary point at $k^{*}=1.2144$ with
$E=67.1408$ J/correct and a gap of $-1.9\times10^{-7}$, which is zero to numerical
precision. Since $1.2144<4$, a positive gap at $k=4$ is exactly what the theory
requires. The H3 test asked whether the gate-selected depth is stationary, and
the answer is no; that is a statement about the selected depth, not a failure of
the identity.

\emph{What the interpolation does and does not establish.} The near-zero gap at
$k^{*}=1.2144$ is close to a consistency check on the solver: any smooth
interpolant, minimised numerically, will report $E'\approx0$ at its own
minimiser, so the zero is not independent evidence for the location. The
independent content is the sign at $k=4$, which is estimated from the data and
has an interval that excludes zero, and the flatness of $E$ between $k=1$ and
$k=2$: the two values are $73.71$ and $73.33$ J/correct, a difference of $0.5\%$,
with overlapping percentile intervals $[68.2,79.9]$ and $[68.1,79.3]$. The
existence of an interior unconstrained optimum is therefore supported on this
arm, while its location inside $[1,2]$ is not identified. The defensible
statement is that $k=4$ is over-retrieval relative to an unconstrained optimum
near $k=1.2$, not that $1.214$ is the optimal depth. Depth $0$ is excluded from
the identity because $\varepsilon$ is degenerate there and because $A(0)=0.001$
rests on a single correct answer.

\subsection{What is not claimed}\label{app:notclaimed}

The appendix is deliberately narrower than a general theory of retrieval depth.
It does not claim the following.

\begin{itemize}
  \item \emph{A universal $k^{*}$.} Every number quoted here is a property of one
  reader, one retriever, one corpus, one grid, and one hardware condition. The
  comparative statics are signs, not magnitudes, and they are derived for the
  unconstrained ratio.
  \item \emph{That accuracy is concave in general.} Proposition~\ref{app:prop:unimodal}
  is conditional. The P4 validation rejects the coverage family on three of four
  arms and finds non-monotone curves on two, so the hypothesis is tested where it
  is used and is not a maintained assumption about RAG.
  \item \emph{That a zero gap certifies optimality.} The identity is necessary,
  not sufficient, as the convex counterexample in Section~\ref{app:lemma1} shows.
  A gap interval that contains zero is consistent with stationarity; it does not
  establish a minimum, and it cannot speak for a boundary depth.
  \item \emph{That the grid argmin is the continuous argmin.} The grid is coarse
  and uneven, the interpolant is a modelling choice, and the policy object is the
  constrained minimiser of (A3). The two optima are reported separately and are
  never substituted for one another.
  \item \emph{That the ceiling of Proposition~\ref{app:prop:price} is attained by
  a deployable policy.} It is a bound with a sharpness example in a stylized
  mixture. A learned policy is a separate object that must be trained, selected,
  and evaluated on disjoint splits, and the P3 result shows that a realizable
  policy can capture none of the gap.
  \item \emph{That $N^{*}$ is a deployment guarantee.} It is a deterministic
  break-even: no discounting, no uncertainty in the saving, no cost of capital,
  and a certificate that is assumed to settle the question once paid.
  \item \emph{Any claim outside the measured condition.} No carbon claim, no
  claim about other hardware, readers, or retrievers, and no claim that the
  exploratory arms enter the theory.
  \item \emph{Any re-adjudication.} The appendix proves statements; it does not
  change the status of H3 or of any gate. H3 remains adjudicated at $k=4$ and
  remains unsupported.
\end{itemize}

%
%  The manuscript preamble loads amsmath, amssymb, and booktabs but not amsthm,
%  so results are set as \paragraph run-in headings and every proof ends with
%  \hfill$\square$. No preamble change is required, and this file can be built
%  on its own inside the same class.
%
%  Scope: this appendix proves only the statements used in the main text.
%  Values marked "P1 diagnostic" are registered outputs of
%  analysis/squad_qwen4b_p1/ and analysis/identity_continuous.json; they are
%  quoted, not re-estimated, and no gate is re-adjudicated here.
% ===========================================================================
\section{Formal theory of the depth decision}\label{app:theory}

This appendix states and proves the six results the manuscript relies on: the
first-order identity (Lemma~1), the shape condition that makes the identity
useful (Proposition~1), the comparative statics of the optimum (Corollary~1),
the boundary condition for a zero optimum (Corollary~2), the ceiling on the value
of per-question allocation (Proposition~2), and the cost of certifying a depth
(Proposition~3). Section~\ref{app:grid} separates the grid identity from the
continuous optimum, and Section~\ref{app:notclaimed} lists what the theory does
not license. Every result is conditional on the assumptions in
Section~\ref{app:assumptions}; none of them is a statement about the reader, the
retriever, or the hardware as such.

\subsection{Setting and assumptions}\label{app:assumptions}

Depth $k$ is a count of passages and ranges over the closed interval $[0,K]$ with
$K=32$. The measured grid is
\begin{equation}\label{app:eq:grid}
  \mathcal{K}=\{0,1,2,4,8,16,32\},
\end{equation}
so the grid is coarse and unevenly spaced, with gaps of $1$, $1$, $2$, $4$, $8$,
and $16$ passages. Mean energy per query is $J(k)$ in joules per query, accuracy
is $A(k)$ as a proportion, and energy per correct answer is the ratio estimand of
the main text,
\begin{equation}\label{app:eq:ratio}
  E(k)=\frac{J(k)}{A(k)}
  \qquad [\mathrm{J/correct}].
\end{equation}
The two elasticities and their gap are used throughout:
\begin{equation}\label{app:eq:elasticities}
  \varepsilon_{J}(k)=\frac{kJ'(k)}{J(k)},
  \qquad
  \varepsilon_{A}(k)=\frac{kA'(k)}{A(k)},
  \qquad
  \Delta(k)=\varepsilon_{J}(k)-\varepsilon_{A}(k),
\end{equation}
both dimensionless for $k>0$. The assumptions are as follows.

\begin{enumerate}
  \item[(A1)] \emph{Affine energy.} $J(k)=a+bk$ for $k\in[0,K]$, with $a>0$ in
  J/query and $b>0$ in J per passage per query. Equivalently
  $\varepsilon_{J}(k)=bk/(a+bk)\in(0,1)$ and $\varepsilon_{J}$ is strictly
  increasing in $k$.

  \item[(A2)] \emph{Positive accuracy.} $A(k)\in(0,1]$ for $k\in[0,K]$, and $A$ is
  twice differentiable wherever a derivative appears. Curvature signs are stated
  per result and are never assumed globally without being named.

  \item[(A3)] \emph{Decision object.} The policy quantity is the constrained
  minimiser over the grid, $k^{*}_{m}=\arg\min_{k\in\mathcal{F}_{m}}E(k)$ with
  $\mathcal{F}_{m}=\{k\in\mathcal{K}:A(k)\ge A(K)-m\}$, which is the object of the
  main text. Statements about a continuous $k^{*}$ are statements about the
  extension in (A4), not about unmeasured depths.

  \item[(A4)] \emph{Continuum extension.} Off the grid, $A$ denotes a twice
  differentiable interpolant of the seven grid means (the registered analysis
  uses a monotone cubic). The interpolant is a modelling device: it makes
  $E'(k)$ well defined, and every off-grid number inherits its smoothness
  assumption.

  \item[(A5)] \emph{Sampling.} Questions are independent draws from a fixed
  distribution, each question is answered at every depth in $\mathcal{K}$, and
  correctness is binary, $C_{qk}\in\{0,1\}$ with $\mathbb{E}[C_{qk}]=A(k)$. This
  assumption is used only in Proposition~\ref{app:prop:cert}(b).

  \item[(A6)] \emph{Calibration.} Calibrated energy is a positive depth-common
  multiple of the measured energy, $J_{\mathrm{cal}}(k)=fJ(k)$ with $f>0$. The
  calibration rule of the main text is all-or-nothing across depths, so this is
  the regime in which absolute energy is reported.
\end{enumerate}

Assumption (A6) has an immediate consequence that justifies using calibrated
energy in the theory: $E_{\mathrm{cal}}=fE$, so the minimiser, every ratio of two
depths, and every elasticity $\varepsilon_{J}$ are unchanged, because $f$ cancels
in $kJ'_{\mathrm{cal}}/J_{\mathrm{cal}}$. A depth-dependent calibration error
would not cancel, which is why the measurement protocol gates on a per-depth
factor and applies it only if all depths pass.

\begin{table}[pos=htbp]
  \centering
  \caption{Summary of the appendix. Each result is conditional on the
  hypotheses named in its statement; none of them is a claim about a particular
  reader, corpus, or GPU.}
  \label{tab:theorymap}
  \footnotesize
  \setlength{\tabcolsep}{4pt}
  \begin{tabular}{@{}>{\raggedright\arraybackslash}p{1.9cm}>{\raggedright\arraybackslash}p{7.4cm}>{\raggedright\arraybackslash}p{5.0cm}@{}}
    \toprule
    \textbf{Result} & \textbf{What is proved} & \textbf{What it licenses} \\
    \midrule
    Lemma~1 & $E'(k)=0$ iff the two elasticities agree, for $k>0$ &
    A local diagnostic; necessary, not sufficient \\
    Proposition~1 & Concave $A$ makes the stationarity function monotone, so $E$
    is unimodal & Existence and uniqueness of an interior optimum, and a test of
    the hypothesis \\
    Corollary~1 & $\partial k^{*}/\partial a>0$ and $\partial k^{*}/\partial b<0$
    at a slack-gate interior optimum & Directional, qualitative predictions only \\
    Corollary~2 & A closed-form $b/a$ threshold for $k^{*}=0$ under coverage &
    Boundary optima are possible; a check on the coverage family only \\
    Proposition~2 & Any allocation at matched accuracy satisfies
    $E_{\pi}/E_{s}\le 1+bK/a$ & A ceiling on the value of adaptivity, with a
    sharpness example \\
    Proposition~3 & Break-even volume and sample complexity of certification &
    A cost, not a guarantee; the accuracy side dominates \\
    \bottomrule
  \end{tabular}
\end{table}

\subsection{Lemma 1: the first-order identity}\label{app:lemma1}

\paragraph{Lemma 1.} Under (A1) and (A2), for every $k\in(0,K]$,
\begin{equation}\label{app:eq:identity}
  E'(k)=\frac{bA(k)-J(k)A'(k)}{A(k)^{2}}
  =\frac{J(k)}{kA(k)}\,\Delta(k),
\end{equation}
and therefore
\begin{equation}\label{app:eq:identity-iff}
  E'(k)=0
  \iff
  bA(k)=J(k)A'(k)
  \iff
  \frac{kJ'(k)}{J(k)}=\frac{kA'(k)}{A(k)}
  \iff
  \Delta(k)=0 .
\end{equation}
At $k=0$ only the unmultiplied form is meaningful,
$E'(0)=[bA(0)-aA'(0)]/A(0)^{2}$, because the elasticity form is identically
$0=0$ there.

\paragraph{Proof.} By the quotient rule and (A1),
\begin{equation}\label{app:eq:proof-quotient}
  E'(k)=\frac{J'(k)A(k)-J(k)A'(k)}{A(k)^{2}}
  =\frac{bA(k)-J(k)A'(k)}{A(k)^{2}},
\end{equation}
which is the first equality in \eqref{app:eq:identity}; $A(k)>0$ by (A2), so the
expression is defined. For the second equality, write
$E=J/A$ and differentiate logarithmically, $E'/E=J'/J-A'/A$. Multiplying by $k$
gives $kE'(k)/E(k)=\Delta(k)$ and hence
\begin{equation}\label{app:eq:proof-second}
  E'(k)=\frac{E(k)}{k}\Delta(k)=\frac{J(k)}{kA(k)}\Delta(k),
\end{equation}
which agrees with \eqref{app:eq:proof-quotient} after multiplying both sides by
$kA(k)^{2}$. Since $J(k)>0$, $A(k)>0$, and $k>0$, the prefactor $J/(kA)$ in
\eqref{app:eq:proof-second} is strictly positive, so $E'(k)=0$ holds if and only
if $\Delta(k)=0$; the middle equivalence in \eqref{app:eq:identity-iff} is the
same statement multiplied by $k>0$ and divided by the positive quantities $J$
and $A$. \hfill$\square$

\paragraph{Necessity, not sufficiency.} An interior local extremum of a
differentiable $E$ must satisfy \eqref{app:eq:identity-iff}, so the identity is a
necessary condition for an interior optimum and it identifies the candidates.
It is not sufficient: a point where the elasticities agree may be a local
maximum, and the identity says nothing about a boundary optimum. The smallest
counterexample uses a convex accuracy curve, $a=b=1$ and
$A(k)=0.1+0.05k^{2}$. Then $\Delta(k)=0$ at $k=\sqrt{3}-1\approx0.732$, where
$E\approx13.66$ J/correct, while $E(0)=E(2)=10$ J/correct; the stationary point
is a local maximum, and the hypothesis $A''\le0$ of Proposition~\ref{app:prop:unimodal}
fails. A zero gap is therefore evidence about the local slope of $E$ at a depth,
never by itself evidence that the depth is optimal.

\paragraph{Reading the sign.} Because the prefactor in \eqref{app:eq:proof-second}
is positive, $\operatorname{sign}E'(k)=\operatorname{sign}\Delta(k)$ for $k>0$. A
positive gap means the marginal joule buys less than the marginal correct answer,
so a small increase in depth raises energy per correct answer; a negative gap
means the opposite. This is the interpretation used for the reported gap, and it
is invariant to the depth-common factor $f$ of (A6), because both elasticities
are logarithmic derivatives.

\subsection{Proposition 1: unimodality under concave accuracy}\label{app:prop:unimodal}

\paragraph{Proposition 1.} Assume (A1), (A2), and $A''(k)\le0$ on $(0,K)$; write
$J(k)=a+bk$ and define the stationarity function
\begin{equation}\label{app:eq:N}
  N(k)=bA(k)-J(k)A'(k)=A(k)^{2}E'(k)
  \qquad [\mathrm{J\,query^{-1}\,passage^{-1}}].
\end{equation}
Then:
\begin{enumerate}
  \item[(i)] $N'(k)=-(a+bk)A''(k)\ge0$, so $N$ is nondecreasing on $[0,K]$, and
  strictly increasing if $A''<0$ on $(0,K)$.
  \item[(ii)] $E'=N/A^{2}$ changes sign at most once on $(0,K)$, and any sign
  change is from negative to positive. Hence $E$ is nonincreasing then
  nondecreasing: it is unimodal (quasiconvex) and has no interior local maximum.
  If $A''<0$, $E$ has at most one interior stationary point, and when it exists it
  is the unique global minimiser of $E$ on $[0,K]$.
  \item[(iii)] If $A''<0$, an interior minimiser exists if and only if
  \begin{equation}\label{app:eq:interior}
    bA(0)<aA'(0)
    \quad\text{and}\quad
    bA(K)>J(K)A'(K);
  \end{equation}
  otherwise the minimiser is $k=0$ (if the first inequality fails) or $k=K$ (if
  the second fails).
\end{enumerate}

\paragraph{Proof.} (i) Differentiate \eqref{app:eq:N} using (A1):
\begin{equation}\label{app:eq:Nprime}
  N'(k)=bA'(k)-\left[bA'(k)+(a+bk)A''(k)\right]=-(a+bk)A''(k),
\end{equation}
and $a+bk>0$ by (A1), so $N'\ge0$; if $A''<0$ then $N'>0$.

(ii) By \eqref{app:eq:N}, $E'=N/A^{2}$ and $A^{2}>0$, so $E'$ has the sign of
$N$. A nondecreasing function crosses zero at most once, and when it does the
crossing is from negative to positive; hence there is a point $\bar k$ (possibly
$0$ or $K$, possibly an interval when $N$ is flat at zero) with $E'\le0$ for
$k<\bar k$ and $E'\ge0$ for $k>\bar k$. That is unimodality, and no interior local
maximum is possible. If $A''<0$ then $N$ is strictly increasing, so its zero set
is empty or a single point $k^{*}$. At such a point,
\begin{equation}\label{app:eq:Esecond}
  E''(k^{*})=\frac{N'(k^{*})A(k^{*})^{2}-2A(k^{*})A'(k^{*})N(k^{*})}{A(k^{*})^{4}}
  =\frac{N'(k^{*})}{A(k^{*})^{2}}>0,
\end{equation}
because $N(k^{*})=0$; the point is a strict local minimum, and by the sign pattern
it is the unique global minimiser on $[0,K]$.

(iii) If $A''<0$ then $N$ is continuous and strictly increasing. An interior zero
exists if and only if $N(0)<0<N(K)$; substituting \eqref{app:eq:N} gives
$bA(0)-aA'(0)<0$ and $bA(K)-(a+bK)A'(K)>0$, which is \eqref{app:eq:interior}.
If $N(0)\ge0$ then $N>0$ on $(0,K]$ and $E$ is nondecreasing, so the minimiser is
$k=0$; if $N(K)\le0$ then $N<0$ on $[0,K)$ and $E$ is nonincreasing, so the
minimiser is $k=K$. \hfill$\square$

\paragraph{Two remarks.} First, the proof never uses $A'\ge0$: monotonicity of
$N$ depends only on the sign of $A''$. The proposition therefore also covers a
concave decline after an accuracy peak, which is the qualitative shape of the two
non-monotone P4 arms (a peak followed by a decline), although concavity of the
decline is not tested there. Second, the hypothesis is testable on the grid: the discrete
analogue of $A''\le0$ is that the secant slopes between adjacent grid depths are
non-increasing, which is what the companion implementation checks. The hypothesis
is not a description of every accuracy curve; where it fails, Proposition~1 gives
no global statement and the identity remains local.

\paragraph{Existence in elasticity form.} Conditions \eqref{app:eq:interior} can
be read as bounds on the elasticities. The first is
$b/a<A'(0)/A(0)$, i.e.\ the relative accuracy gain per passage at zero must exceed
the relative energy cost per passage. The second is
$\varepsilon_{A}(K)<\varepsilon_{J}(K)=bK/(a+bK)<1$, i.e.\ the accuracy
elasticity at the top of the grid must be below the energy elasticity. Both are
properties of the fitted curves, not of the decision rule, and both are evaluated
without the accuracy gate of (A3).

\subsection{Corollary 1: comparative statics of the optimum}\label{app:cor:statics}

\paragraph{Corollary 1.} Assume (A1), (A2), $A'(k^{*})>0$ and $A''(k^{*})<0$ at an
interior stationary point $k^{*}=k^{*}(a,b)$ of $E$. Then $k^{*}$ is continuously
differentiable in $(a,b)$ in a neighbourhood, and
\begin{equation}\label{app:eq:statics}
  \frac{\partial k^{*}}{\partial a}
  =\frac{-A'(k^{*})}{(a+bk^{*})A''(k^{*})}>0,
  \qquad
  \frac{\partial k^{*}}{\partial b}
  =\frac{A(k^{*})-k^{*}A'(k^{*})}{(a+bk^{*})A''(k^{*})}
  =-\frac{aA(k^{*})}{(a+bk^{*})^{2}\left(-A''(k^{*})\right)}<0 .
\end{equation}

\paragraph{Proof.} By Lemma~1 the stationary point solves
$F(k,a,b)=bA(k)-(a+bk)A'(k)=0$, and $F$ is continuously differentiable under
(A2). Its partials are
\begin{equation}\label{app:eq:Fpartials}
  F_{k}=-(a+bk)A''(k),
  \qquad
  F_{a}=-A'(k),
  \qquad
  F_{b}=A(k)-kA'(k).
\end{equation}
At $k^{*}$, $F_{k}=-(a+bk^{*})A''(k^{*})>0$ because $A''(k^{*})<0$; this is also
the second-order condition \eqref{app:eq:Esecond}. The implicit function theorem
therefore applies, and
\begin{equation}\label{app:eq:IFT}
  \frac{\partial k^{*}}{\partial a}=-\frac{F_{a}}{F_{k}},
  \qquad
  \frac{\partial k^{*}}{\partial b}=-\frac{F_{b}}{F_{k}} .
\end{equation}
Substituting \eqref{app:eq:Fpartials} gives the first equality in each line of
\eqref{app:eq:statics}. The first sign is immediate since $A'(k^{*})>0$ and the
denominator is negative. For the second, apply the identity
\eqref{app:eq:identity-iff} at $k^{*}$:
$k^{*}A'(k^{*})/A(k^{*})=\varepsilon_{J}(k^{*})=bk^{*}/(a+bk^{*})<1$, hence
\begin{equation}\label{app:eq:Fbpositive}
  A(k^{*})-k^{*}A'(k^{*})
  =A(k^{*})\left(1-\varepsilon_{J}(k^{*})\right)
  =\frac{aA(k^{*})}{a+bk^{*}}>0,
\end{equation}
and the denominator of the second line is negative, so
$\partial k^{*}/\partial b<0$. \hfill$\square$

\paragraph{Reading.} A larger fixed cost per query makes depth relatively cheap
and pushes the unconstrained optimum deeper; a larger per-passage cost makes
depth relatively expensive and pushes it shallower. The corollary is local and
applies only when the accuracy gate is slack: when the gate binds, the policy
depth is the smallest feasible grid depth and is locally insensitive to $(a,b)$,
so the derivative is zero in the one-sided sense at that corner. It is also
silent about the magnitude, because $A'(k^{*})$ and $A''(k^{*})$ are estimated
from the same finite sample as $A(k^{*})$; the sign is a prediction, the size is
not.

\subsection{Corollary 2: when the optimum is the boundary}\label{app:cor:boundary}

\paragraph{Corollary 2.} Assume (A1) and the stylized coverage family
\begin{equation}\label{app:eq:coverage}
  A(k)=1-(1-a_{0})(1-p)^{k},
  \qquad
  a_{0}=A(0)\in(0,1),
  \qquad
  p\in(0,1),
\end{equation}
so that $A(k)\in(a_{0},1)$ and (A2) holds on $[0,\infty)$. Let
$L=-\ln(1-p)>0$. Then $A'(k)=L\left(1-A(k)\right)>0$ and
$A''(k)=-L^{2}\left(1-A(k)\right)<0$, Proposition~\ref{app:prop:unimodal}
applies, and
\begin{equation}\label{app:eq:boundary}
  k^{*}=0
  \iff
  \frac{b}{a}\ge\frac{(1-a_{0})\left(-\ln(1-p)\right)}{a_{0}}
  =\frac{1-a_{0}}{a_{0}}\,L .
\end{equation}
When the inequality is reversed, the unique interior optimum solves
\begin{equation}\label{app:eq:coverage-optimum}
  (1-a_{0})(1-p)^{k^{*}}=\frac{b}{aL+b+bLk^{*}} .
\end{equation}

\paragraph{Proof.} Differentiating \eqref{app:eq:coverage} gives
$A'(k)=-(1-a_{0})(1-p)^{k}\ln(1-p)=L\left(1-A(k)\right)$ and
$A''(k)=-L^{2}\left(1-A(k)\right)<0$, so $A$ is strictly increasing and strictly
concave and (A2) holds. Evaluate \eqref{app:eq:N} at zero:
\begin{equation}\label{app:eq:N0}
  N(0)=bA(0)-aA'(0)=ba_{0}-a(1-a_{0})L
  =a_{0}\left[b-a\frac{1-a_{0}}{a_{0}}L\right].
\end{equation}
By Proposition~\ref{app:prop:unimodal}(i), $N$ is strictly increasing. As
$k\to\infty$, $A(k)\to1$ and $(a+bk)(1-A(k))\to0$ because the exponential decay
dominates the linear factor, so
\begin{equation}\label{app:eq:Nlimit}
  N(k)=bA(k)-(a+bk)L\left(1-A(k)\right)\longrightarrow b>0 .
\end{equation}
If $N(0)\ge0$, then $N>0$ on $(0,\infty)$ and $E$ is strictly increasing, so
$k^{*}=0$; by \eqref{app:eq:N0} this is exactly the inequality in
\eqref{app:eq:boundary}. If $N(0)<0$, then \eqref{app:eq:Nlimit} and strict
monotonicity give a unique zero $k^{*}\in(0,\infty)$, which is the unique
minimiser by Proposition~\ref{app:prop:unimodal}(ii). Setting $N(k^{*})=0$ and
writing $s=1-A(k^{*})=(1-a_{0})(1-p)^{k^{*}}$ gives
$b(1-s)=(a+bk^{*})Ls$, i.e.\ $b=s\left(aL+b+bLk^{*}\right)$, which is
\eqref{app:eq:coverage-optimum}. \hfill$\square$

\paragraph{Scope.} The corollary is a statement about one stylized family, not
about the observed curves. The P4 validation rejects the family as a general
description of the four arms ($R^{2}$ between $0.097$ and $0.879$), so its role
is to show that a boundary optimum is a possible outcome of diminishing returns,
and to give the threshold in closed form. It does not certify any observed
$k^{*}$, and it is not used to choose a depth.

\subsection{Proposition 2: the price of static depth}\label{app:prop:price}

\paragraph{Proposition 2.} Assume (A1). Let $\pi$ be any per-question depth
assignment with values in $[0,K]$, with population energy
$J_{\pi}=\mathbb{E}_{q}[J(\pi(q))]$ and accuracy
$A_{\pi}=\mathbb{E}_{q}[A(\pi(q))]$, and let $k_{s}$ be a static depth with
accuracy $A_{s}=A(k_{s})$. If $A_{\pi}\ge A_{s}$, then
\begin{equation}\label{app:eq:ceiling}
  \frac{E_{\pi}}{E_{s}}
  =\frac{J_{\pi}}{J(k_{s})}\cdot\frac{A_{s}}{A_{\pi}}
  \le\frac{a+bK}{a}
  =\frac{1}{1-\varepsilon_{J}(K)} ,
\end{equation}
and the largest saving any such assignment can deliver is the variable energy
share at the reference depth,
\begin{equation}\label{app:eq:ceiling-saving}
  1-\frac{E_{\pi}}{E_{s}}\le\frac{bK}{a+bK}=\varepsilon_{J}(K)<1 .
\end{equation}
The bound is sharp. For $\varepsilon\in(0,1)$ let a fraction $1-\varepsilon$ of
questions be easy, answered correctly at every depth including $k=0$, and a
fraction $\varepsilon$ be hard, answered correctly only at depths $k\ge K$. Then
the best static depth at accuracy $1$ is $k_{s}=K$, the hindsight oracle
$\pi_{\mathrm{or}}$ assigns easy questions to $0$ and hard questions to $K$, and
\begin{equation}\label{app:eq:rho}
  \rho(\varepsilon)
  =\frac{E_{\mathrm{static}}}{E_{\mathrm{oracle}}}
  =\frac{a+bK}{a+\varepsilon bK}
  \;\longrightarrow\;
  1+\frac{bK}{a}
  =\frac{1}{1-\varepsilon_{J}(K)}
  \qquad(\varepsilon\to0).
\end{equation}
The ratio is written here as static over oracle, so that the quantity is the
\emph{price} of fixing one depth; written in the opposite order (oracle over
static) the same limit is its reciprocal $a/(a+bK)<1$.

\paragraph{Proof.} By (A1), $J(\pi(q))\in[a,a+bK]$ for every question, so
$a\le J_{\pi}\le a+bK$ and $J(k_{s})\ge a$; hence
$J_{\pi}/J(k_{s})\le(a+bK)/a$. Since $A_{\pi}\ge A_{s}>0$, the factor
$A_{s}/A_{\pi}$ is at most one, which proves \eqref{app:eq:ceiling}. For
\eqref{app:eq:ceiling-saving}, subtract the reciprocal:
$1-a/(a+bK)=bK/(a+bK)=\varepsilon_{J}(K)$ by (A1).

For sharpness, let $q$ be easy with probability $1-\varepsilon$ and hard with
probability $\varepsilon$. Every static depth $k<K$ answers only the easy
questions, so its accuracy is $1-\varepsilon<1$; the only static depths with
accuracy $1$ are $k\ge K$, and among them $J$ is minimised at $k=K$. Hence
$E_{\mathrm{static}}=a+bK$. The oracle answers every question correctly, so its
accuracy is $1$ and its energy is
$(1-\varepsilon)a+\varepsilon(a+bK)=a+\varepsilon bK$, giving
$E_{\mathrm{oracle}}=a+\varepsilon bK$ and the first equality in
\eqref{app:eq:rho}. The limit is immediate, and the last equality repeats
$bK/(a+bK)=\varepsilon_{J}(K)$. \hfill$\square$

\paragraph{Scope.} Proposition~\ref{app:prop:price} is an upper bound with a
sharpness example, not a lower bound on the value of adaptivity and not a claim
about the observed workload. Two consequences matter for the paper. First, the
ceiling is a property of the affine energy model alone: it caps \emph{any}
allocation, learned or oracle, at matched accuracy, and it is the formal reason
the oracle gain is reported as a fraction of the static energy rather than as an
unbounded promise. Second, the example shows that the ceiling is approached when
a small hard stratum forces the whole workload to the deepest setting; when the
stratum is large ($\varepsilon\to1$), the ratio tends to one and adaptivity is
worth nothing. The empirical matched-oracle gain reported for P3 is far below the
P1 ceiling of $\varepsilon_{J}(32)=0.7517$, as it must be, and the pilot oracle
figure of $78.8\%$ is below the pilot ceiling
$bK/(a+bK)=0.8227$ computed from $a=15.74$ and $b=2.283$; both are consistency
checks on registered quantities, not new evidence.

\subsection{Proposition 3: the cost of certification}\label{app:prop:cert}

Certification has two separate costs: the energy of the experiment (part a) and
the number of questions needed for the statistical statement (part b). They are
stated separately because the first is a deterministic break-even and the second
is a sampling requirement, and only their product is the cost of a certified
depth.

\paragraph{Proposition 3(a): economic break-even.} Assume (A1), a deployment of
$N$ queries drawn from the same distribution as the experiment, one depth used
for all $N$ queries, no discounting, and a one-off certification cost
$C_{\mathrm{cert}}>0$ in joules. Let $k_{0}=K$ be the reference depth and
$k^{*}$ the candidate. Define the per-query break-even and the
per-correct-answer break-even by
\begin{equation}\label{app:eq:nstar}
  N^{*}_{J}=\frac{C_{\mathrm{cert}}}{J(k_{0})-J(k^{*})},
  \qquad
  N^{*}_{A}=\frac{C_{\mathrm{cert}}}{A(k^{*})\left[E(k_{0})-E(k^{*})\right]},
\end{equation}
both in queries and both defined only when the denominator is positive. Then
\begin{equation}\label{app:eq:nstar-identity}
  A(k^{*})\left[E(k_{0})-E(k^{*})\right]
  =\frac{A(k^{*})}{A(k_{0})}J(k_{0})-J(k^{*}),
\end{equation}
so $N^{*}_{A}=N^{*}_{J}$ if and only if $A(k^{*})=A(k_{0})$, and otherwise
\begin{equation}\label{app:eq:beta}
  \beta:=\frac{N^{*}_{A}}{N^{*}_{J}}
  =\frac{J(k_{0})-J(k^{*})}{\dfrac{A(k^{*})}{A(k_{0})}J(k_{0})-J(k^{*})},
\end{equation}
with $\beta\ge1$ if and only if $A(k^{*})\le A(k_{0})$.

\paragraph{Proof.} Equation \eqref{app:eq:nstar-identity} is
$A(k^{*})E(k_{0})-A(k^{*})E(k^{*})=A(k^{*})J(k_{0})/A(k_{0})-J(k^{*})$, using
$A(k^{*})E(k^{*})=J(k^{*})$. Dividing the two lines of \eqref{app:eq:nstar}
gives \eqref{app:eq:beta}. For the sign, compare denominators:
\begin{equation}\label{app:eq:beta-sign}
  \left[J(k_{0})-J(k^{*})\right]
  -\left[\frac{A(k^{*})}{A(k_{0})}J(k_{0})-J(k^{*})\right]
  =J(k_{0})\left(1-\frac{A(k^{*})}{A(k_{0})}\right),
\end{equation}
which is non-negative exactly when $A(k^{*})\le A(k_{0})$. Both denominators are
positive by hypothesis, so the inequality passes to $\beta$.
\hfill$\square$

\paragraph{Interpretation and the approximation.} $N^{*}_{J}$ is the number of
deployment queries at which the energy saved by switching depth equals the
certificate: after $N$ queries at $k^{*}$ instead of $k_{0}$ the saving is
$N\left[J(k_{0})-J(k^{*})\right]$, and it repays $C_{\mathrm{cert}}$ at
$N=N^{*}_{J}$. $N^{*}_{A}$ answers a different question: it values the candidate's
correct-answer output at the per-correct-answer saving $E(k_{0})-E(k^{*})$, so $N$
candidate-depth queries produce $N A(k^{*})$ correct answers and a cumulative
saving $N A(k^{*})\left[E(k_{0})-E(k^{*})\right]$, which equals
$C_{\mathrm{cert}}$ at $N=N^{*}_{A}$. The two coincide only at
equal accuracy. When $A(k^{*})<A(k_{0})$, the paper's per-correct-answer formula
is therefore an approximation to the deployment-queries reading, and it errs on
the conservative side ($\beta>1$): it reports a longer break-even because it also
requires the candidate to reproduce the correct-answer output that its lower
accuracy would otherwise lose. The direction is known but the formula is not
exact, which is why both numbers are reported rather than one.

\paragraph{P1 diagnostic.} With the registered P1 values
$C_{\mathrm{cert}}=311{,}992.064$ J, $J(4)=36.2537$ J/query,
$J(32)=94.9450$ J/query, $A(4)=0.450$, and $A(32)=0.493$, one obtains
$N^{*}_{J}=5{,}315.8$ queries, $N^{*}_{A}=6{,}189.1$ queries, and
$\beta=1.1643$. The certification cost is itself a sum over the grid,
$C_{\mathrm{cert}}=n\sum_{k\in\mathcal{K}}J(k)=1{,}000\times311.992$ J, so the
sweep energy and the reported break-even use the same accounting.

\paragraph{Proposition 3(b): statistical sample complexity.} Assume (A5), let
$m$ be the accuracy margin (written $\delta$ in the main text, and renamed here
to avoid a clash with the error probability), and let
$\Delta=A(k_{0})-A(k^{*})>0$. For the paired differences
$D_{q}=C_{qk_{0}}-C_{qk^{*}}\in\{-1,0,1\}$ with
$\mathbb{E}[D_{q}]=\Delta$, put
$t_{n}=\sqrt{2\ln(1/\alpha)/n}$ and declare non-inferiority of $k^{*}$ when
\begin{equation}\label{app:eq:declare}
  \bar D+t_{n}<m .
\end{equation}
Then:
\begin{enumerate}
  \item[(i)] \eqref{app:eq:declare} is a level-$\alpha$ test of
  $H_{0}:\Delta\ge m$.
  \item[(ii)] At a true difference $\Delta=m_{1}<m$ the probability of failing to
  declare non-inferiority is at most
  \begin{equation}\label{app:eq:hoeffding}
    \Pr\left(\text{fail to declare}\mid \Delta=m_{1}\right)
    \le\exp\left(-\frac{n\left(\left(m-m_{1}\right)-t_{n}\right)^{2}}{2}\right),
  \end{equation}
  so $n\ge8\max\left(\ln(1/\alpha),\ln(1/\delta_{\mathrm{err}})\right)/(m-m_{1})^{2}$
  paired questions suffice for a failure probability at most
  $\delta_{\mathrm{err}}$.
  \item[(iii)] Selecting the more accurate of two depths at a true gap $\Delta$
  needs $n\ge2\ln(1/\delta_{\mathrm{err}})/\Delta^{2}$, and selecting among $d$
  depths whose smallest gap to the optimum is $\Delta_{\min}$ needs
  $n\ge2\ln\left((d-1)/\delta_{\mathrm{err}}\right)/\Delta_{\min}^{2}$ by a union
  bound; in both cases the correct depth is selected with probability at least
  $1-\delta_{\mathrm{err}}$.
\end{enumerate}

\paragraph{Proof sketch.} $D_{q}$ is bounded in $[-1,1]$ and has range $2$, so
Hoeffding's inequality gives, for every $u>0$,
\begin{equation}\label{app:eq:two-tails}
  \Pr\left(\bar D-\Delta\le-u\right)\le\exp\left(-\frac{nu^{2}}{2}\right),
  \qquad
  \Pr\left(\bar D-\Delta\ge u\right)\le\exp\left(-\frac{nu^{2}}{2}\right).
\end{equation}
(i) Under $H_{0}$, $\Delta\ge m$, so declaring non-inferiority implies
$\bar D-\Delta<-t_{n}$, which has probability at most
$\exp(-nt_{n}^{2}/2)=\alpha$.

(ii) Failure to declare means $\bar D+t_{n}\ge m$, hence
$\bar D-\Delta\ge(m-m_{1})-t_{n}$, and the upper tail of
\eqref{app:eq:two-tails} applies. If $t_{n}\le(m-m_{1})/2$, which holds once
$n\ge8\ln(1/\alpha)/(m-m_{1})^{2}$, the exponent is at most
$-n(m-m_{1})^{2}/8$, and the stated $n$ makes it at most
$\ln(1/\delta_{\mathrm{err}})$.

(iii) The selection error is
$\Pr(\bar D\le0)=\Pr(\bar D-\Delta\le-\Delta)\le\exp(-n\Delta^{2}/2)$, and the
union bound over the $d-1$ alternatives gives the second statement.
\hfill$\square$

\paragraph{Scope and consequences.} The bound is distribution-free and makes no
normal approximation, but it is deliberately conservative: it ignores the
pairing, i.e.\ the discordance rate, which is what a McNemar calculation uses.
The tighter normal-approximation requirement scales like
$\psi/(m-m_{1})^{2}$ at discordance $\psi$, so at a fixed difference the observed
P1 discordance of $13.1\%$ requires roughly $1.6$ times the pairs that the
planned $8\%$ implied, and the registered frontier puts the 5 pp requirement at
$n=7{,}138$ pairs. Three consequences are used in the main text. First, the
sample complexity diverges as $m_{1}\uparrow m$: a margin that the data show is
several percentage points wide cannot be resolved by a design powered at a
near-zero difference, which is the formal content of the P1 precision failure.
Second, the energy side does not drive the requirement: per-query energy has a P1
repeat CV of $2.81\%$ at $k=8$, so the relative standard error of a mean energy
at $n=1{,}000$ is about $2.81\%/\sqrt{1{,}000}<0.1\%$, while the accuracy
contrast at 5 pp remains unresolved at the same $n$. The accuracy side is the
binding one. Third, the certification cost is linear in the statistical
requirement, because a sweep of $n$ questions costs
$C_{\mathrm{cert}}(n)=n\,c_{\mathrm{sweep}}$ with
$c_{\mathrm{sweep}}=\sum_{k\in\mathcal{K}}J(k)$ in J/query, so
\begin{equation}\label{app:eq:nstar-n}
  N^{*}_{J}(n)=\frac{n\,c_{\mathrm{sweep}}}{J(k_{0})-J(k^{*})},
\end{equation}
which grows linearly in $n$. The statistical requirement and the economic
break-even therefore cannot be chosen independently, and a depth that needs a
large experiment needs a proportionally larger deployment to justify it.

\subsection{Grid resolution and the H3 diagnostic}\label{app:grid}

The identity of Lemma~1 is a statement about a derivative. On the grid two
approximations intervene, and separating them is what makes the H3 result
readable.

\emph{First, the derivative is a secant.} The reported accuracy elasticity at a
grid depth is a finite difference over the neighbouring depths, so the identity
at a grid point is the statement that the secant elasticity equals the affine
elasticity. At $k=4$ the central difference uses $k=2$ and $k=8$, a step of three
passages; at $k=8$ it uses $k=4$ and $k=16$, a step of six; at the edges the
difference is one-sided, with step one at $k=1$ and step sixteen at $k=32$. A
central difference approximates $A'(k)$ to second order in the step and a
one-sided difference to first order, so a grid gap is a discretisation of
$\Delta(k)$, not a different quantity.

\emph{Second, the decision object is constrained.} The policy depth of (A3) is the
energy minimiser inside the feasible set $\mathcal{F}_{m}$, whereas $k^{*}$ in
Lemma~1 is the unconstrained stationary point of $E$. The two objects coincide
only when the accuracy gate is slack. In P1 the gate binds: the 5 pp rule selects
$k=4$ because it is the smallest feasible depth, not because $E$ is stationary
there. Testing the identity at the gate-selected depth is therefore a test of
local stationarity at a point that the theory does not expect to be stationary,
and the main text reports it as such.

\emph{The P1 numbers.} With $a=24.0759$ J/query and
$b=2.2772$ J per passage per query, the affine elasticity is
$\varepsilon_{J}(4)=0.2745$, the secant elasticity is
$\varepsilon_{A}(4)=0.0741$, and the gap is $\Delta(4)=+0.2004$ with a 95\%
interval of $[0.1677,0.2320]$ that excludes zero. The positive sign says that
$E$ is increasing at $k=4$: the last passages bought at that depth cost more
energy per correct answer than they return. The monotone-cubic interpolant of the
same seven means has its unconstrained stationary point at $k^{*}=1.2144$ with
$E=67.1408$ J/correct and a gap of $-1.9\times10^{-7}$, which is zero to numerical
precision. Since $1.2144<4$, a positive gap at $k=4$ is exactly what the theory
requires. The H3 test asked whether the gate-selected depth is stationary, and
the answer is no; that is a statement about the selected depth, not a failure of
the identity.

\emph{What the interpolation does and does not establish.} The near-zero gap at
$k^{*}=1.2144$ is close to a consistency check on the solver: any smooth
interpolant, minimised numerically, will report $E'\approx0$ at its own
minimiser, so the zero is not independent evidence for the location. The
independent content is the sign at $k=4$, which is estimated from the data and
has an interval that excludes zero, and the flatness of $E$ between $k=1$ and
$k=2$: the two values are $73.71$ and $73.33$ J/correct, a difference of $0.5\%$,
with overlapping percentile intervals $[68.2,79.9]$ and $[68.1,79.3]$. The
existence of an interior unconstrained optimum is therefore supported on this
arm, while its location inside $[1,2]$ is not identified. The defensible
statement is that $k=4$ is over-retrieval relative to an unconstrained optimum
near $k=1.2$, not that $1.214$ is the optimal depth. Depth $0$ is excluded from
the identity because $\varepsilon$ is degenerate there and because $A(0)=0.001$
rests on a single correct answer.

\subsection{What is not claimed}\label{app:notclaimed}

The appendix is deliberately narrower than a general theory of retrieval depth.
It does not claim the following.

\begin{itemize}
  \item \emph{A universal $k^{*}$.} Every number quoted here is a property of one
  reader, one retriever, one corpus, one grid, and one hardware condition. The
  comparative statics are signs, not magnitudes, and they are derived for the
  unconstrained ratio.
  \item \emph{That accuracy is concave in general.} Proposition~\ref{app:prop:unimodal}
  is conditional. The P4 validation rejects the coverage family on three of four
  arms and finds non-monotone curves on two, so the hypothesis is tested where it
  is used and is not a maintained assumption about RAG.
  \item \emph{That a zero gap certifies optimality.} The identity is necessary,
  not sufficient, as the convex counterexample in Section~\ref{app:lemma1} shows.
  A gap interval that contains zero is consistent with stationarity; it does not
  establish a minimum, and it cannot speak for a boundary depth.
  \item \emph{That the grid argmin is the continuous argmin.} The grid is coarse
  and uneven, the interpolant is a modelling choice, and the policy object is the
  constrained minimiser of (A3). The two optima are reported separately and are
  never substituted for one another.
  \item \emph{That the ceiling of Proposition~\ref{app:prop:price} is attained by
  a deployable policy.} It is a bound with a sharpness example in a stylized
  mixture. A learned policy is a separate object that must be trained, selected,
  and evaluated on disjoint splits, and the P3 result shows that a realizable
  policy can capture none of the gap.
  \item \emph{That $N^{*}$ is a deployment guarantee.} It is a deterministic
  break-even: no discounting, no uncertainty in the saving, no cost of capital,
  and a certificate that is assumed to settle the question once paid.
  \item \emph{Any claim outside the measured condition.} No carbon claim, no
  claim about other hardware, readers, or retrievers, and no claim that the
  exploratory arms enter the theory.
  \item \emph{Any re-adjudication.} The appendix proves statements; it does not
  change the status of H3 or of any gate. H3 remains adjudicated at $k=4$ and
  remains unsupported.
\end{itemize}

%
%  The manuscript preamble loads amsmath, amssymb, and booktabs but not amsthm,
%  so results are set as \paragraph run-in headings and every proof ends with
%  \hfill$\square$. No preamble change is required, and this file can be built
%  on its own inside the same class.
%
%  Scope: this appendix proves only the statements used in the main text.
%  Values marked "P1 diagnostic" are registered outputs of
%  analysis/squad_qwen4b_p1/ and analysis/identity_continuous.json; they are
%  quoted, not re-estimated, and no gate is re-adjudicated here.
% ===========================================================================
\section{Formal theory of the depth decision}\label{app:theory}

This appendix states and proves the six results the manuscript relies on: the
first-order identity (Lemma~1), the shape condition that makes the identity
useful (Proposition~1), the comparative statics of the optimum (Corollary~1),
the boundary condition for a zero optimum (Corollary~2), the ceiling on the value
of per-question allocation (Proposition~2), and the cost of certifying a depth
(Proposition~3). Section~\ref{app:grid} separates the grid identity from the
continuous optimum, and Section~\ref{app:notclaimed} lists what the theory does
not license. Every result is conditional on the assumptions in
Section~\ref{app:assumptions}; none of them is a statement about the reader, the
retriever, or the hardware as such.

\subsection{Setting and assumptions}\label{app:assumptions}

Depth $k$ is a count of passages and ranges over the closed interval $[0,K]$ with
$K=32$. The measured grid is
\begin{equation}\label{app:eq:grid}
  \mathcal{K}=\{0,1,2,4,8,16,32\},
\end{equation}
so the grid is coarse and unevenly spaced, with gaps of $1$, $1$, $2$, $4$, $8$,
and $16$ passages. Mean energy per query is $J(k)$ in joules per query, accuracy
is $A(k)$ as a proportion, and energy per correct answer is the ratio estimand of
the main text,
\begin{equation}\label{app:eq:ratio}
  E(k)=\frac{J(k)}{A(k)}
  \qquad [\mathrm{J/correct}].
\end{equation}
The two elasticities and their gap are used throughout:
\begin{equation}\label{app:eq:elasticities}
  \varepsilon_{J}(k)=\frac{kJ'(k)}{J(k)},
  \qquad
  \varepsilon_{A}(k)=\frac{kA'(k)}{A(k)},
  \qquad
  \Delta(k)=\varepsilon_{J}(k)-\varepsilon_{A}(k),
\end{equation}
both dimensionless for $k>0$. The assumptions are as follows.

\begin{enumerate}
  \item[(A1)] \emph{Affine energy.} $J(k)=a+bk$ for $k\in[0,K]$, with $a>0$ in
  J/query and $b>0$ in J per passage per query. Equivalently
  $\varepsilon_{J}(k)=bk/(a+bk)\in(0,1)$ and $\varepsilon_{J}$ is strictly
  increasing in $k$.

  \item[(A2)] \emph{Positive accuracy.} $A(k)\in(0,1]$ for $k\in[0,K]$, and $A$ is
  twice differentiable wherever a derivative appears. Curvature signs are stated
  per result and are never assumed globally without being named.

  \item[(A3)] \emph{Decision object.} The policy quantity is the constrained
  minimiser over the grid, $k^{*}_{m}=\arg\min_{k\in\mathcal{F}_{m}}E(k)$ with
  $\mathcal{F}_{m}=\{k\in\mathcal{K}:A(k)\ge A(K)-m\}$, which is the object of the
  main text. Statements about a continuous $k^{*}$ are statements about the
  extension in (A4), not about unmeasured depths.

  \item[(A4)] \emph{Continuum extension.} Off the grid, $A$ denotes a twice
  differentiable interpolant of the seven grid means (the registered analysis
  uses a monotone cubic). The interpolant is a modelling device: it makes
  $E'(k)$ well defined, and every off-grid number inherits its smoothness
  assumption.

  \item[(A5)] \emph{Sampling.} Questions are independent draws from a fixed
  distribution, each question is answered at every depth in $\mathcal{K}$, and
  correctness is binary, $C_{qk}\in\{0,1\}$ with $\mathbb{E}[C_{qk}]=A(k)$. This
  assumption is used only in Proposition~\ref{app:prop:cert}(b).

  \item[(A6)] \emph{Calibration.} Calibrated energy is a positive depth-common
  multiple of the measured energy, $J_{\mathrm{cal}}(k)=fJ(k)$ with $f>0$. The
  calibration rule of the main text is all-or-nothing across depths, so this is
  the regime in which absolute energy is reported.
\end{enumerate}

Assumption (A6) has an immediate consequence that justifies using calibrated
energy in the theory: $E_{\mathrm{cal}}=fE$, so the minimiser, every ratio of two
depths, and every elasticity $\varepsilon_{J}$ are unchanged, because $f$ cancels
in $kJ'_{\mathrm{cal}}/J_{\mathrm{cal}}$. A depth-dependent calibration error
would not cancel, which is why the measurement protocol gates on a per-depth
factor and applies it only if all depths pass.

\begin{table}[pos=htbp]
  \centering
  \caption{Summary of the appendix. Each result is conditional on the
  hypotheses named in its statement; none of them is a claim about a particular
  reader, corpus, or GPU.}
  \label{tab:theorymap}
  \footnotesize
  \setlength{\tabcolsep}{4pt}
  \begin{tabular}{@{}>{\raggedright\arraybackslash}p{1.9cm}>{\raggedright\arraybackslash}p{7.4cm}>{\raggedright\arraybackslash}p{5.0cm}@{}}
    \toprule
    \textbf{Result} & \textbf{What is proved} & \textbf{What it licenses} \\
    \midrule
    Lemma~1 & $E'(k)=0$ iff the two elasticities agree, for $k>0$ &
    A local diagnostic; necessary, not sufficient \\
    Proposition~1 & Concave $A$ makes the stationarity function monotone, so $E$
    is unimodal & Existence and uniqueness of an interior optimum, and a test of
    the hypothesis \\
    Corollary~1 & $\partial k^{*}/\partial a>0$ and $\partial k^{*}/\partial b<0$
    at a slack-gate interior optimum & Directional, qualitative predictions only \\
    Corollary~2 & A closed-form $b/a$ threshold for $k^{*}=0$ under coverage &
    Boundary optima are possible; a check on the coverage family only \\
    Proposition~2 & Any allocation at matched accuracy satisfies
    $E_{\pi}/E_{s}\le 1+bK/a$ & A ceiling on the value of adaptivity, with a
    sharpness example \\
    Proposition~3 & Break-even volume and sample complexity of certification &
    A cost, not a guarantee; the accuracy side dominates \\
    \bottomrule
  \end{tabular}
\end{table}

\subsection{Lemma 1: the first-order identity}\label{app:lemma1}

\paragraph{Lemma 1.} Under (A1) and (A2), for every $k\in(0,K]$,
\begin{equation}\label{app:eq:identity}
  E'(k)=\frac{bA(k)-J(k)A'(k)}{A(k)^{2}}
  =\frac{J(k)}{kA(k)}\,\Delta(k),
\end{equation}
and therefore
\begin{equation}\label{app:eq:identity-iff}
  E'(k)=0
  \iff
  bA(k)=J(k)A'(k)
  \iff
  \frac{kJ'(k)}{J(k)}=\frac{kA'(k)}{A(k)}
  \iff
  \Delta(k)=0 .
\end{equation}
At $k=0$ only the unmultiplied form is meaningful,
$E'(0)=[bA(0)-aA'(0)]/A(0)^{2}$, because the elasticity form is identically
$0=0$ there.

\paragraph{Proof.} By the quotient rule and (A1),
\begin{equation}\label{app:eq:proof-quotient}
  E'(k)=\frac{J'(k)A(k)-J(k)A'(k)}{A(k)^{2}}
  =\frac{bA(k)-J(k)A'(k)}{A(k)^{2}},
\end{equation}
which is the first equality in \eqref{app:eq:identity}; $A(k)>0$ by (A2), so the
expression is defined. For the second equality, write
$E=J/A$ and differentiate logarithmically, $E'/E=J'/J-A'/A$. Multiplying by $k$
gives $kE'(k)/E(k)=\Delta(k)$ and hence
\begin{equation}\label{app:eq:proof-second}
  E'(k)=\frac{E(k)}{k}\Delta(k)=\frac{J(k)}{kA(k)}\Delta(k),
\end{equation}
which agrees with \eqref{app:eq:proof-quotient} after multiplying both sides by
$kA(k)^{2}$. Since $J(k)>0$, $A(k)>0$, and $k>0$, the prefactor $J/(kA)$ in
\eqref{app:eq:proof-second} is strictly positive, so $E'(k)=0$ holds if and only
if $\Delta(k)=0$; the middle equivalence in \eqref{app:eq:identity-iff} is the
same statement multiplied by $k>0$ and divided by the positive quantities $J$
and $A$. \hfill$\square$

\paragraph{Necessity, not sufficiency.} An interior local extremum of a
differentiable $E$ must satisfy \eqref{app:eq:identity-iff}, so the identity is a
necessary condition for an interior optimum and it identifies the candidates.
It is not sufficient: a point where the elasticities agree may be a local
maximum, and the identity says nothing about a boundary optimum. The smallest
counterexample uses a convex accuracy curve, $a=b=1$ and
$A(k)=0.1+0.05k^{2}$. Then $\Delta(k)=0$ at $k=\sqrt{3}-1\approx0.732$, where
$E\approx13.66$ J/correct, while $E(0)=E(2)=10$ J/correct; the stationary point
is a local maximum, and the hypothesis $A''\le0$ of Proposition~\ref{app:prop:unimodal}
fails. A zero gap is therefore evidence about the local slope of $E$ at a depth,
never by itself evidence that the depth is optimal.

\paragraph{Reading the sign.} Because the prefactor in \eqref{app:eq:proof-second}
is positive, $\operatorname{sign}E'(k)=\operatorname{sign}\Delta(k)$ for $k>0$. A
positive gap means the marginal joule buys less than the marginal correct answer,
so a small increase in depth raises energy per correct answer; a negative gap
means the opposite. This is the interpretation used for the reported gap, and it
is invariant to the depth-common factor $f$ of (A6), because both elasticities
are logarithmic derivatives.

\subsection{Proposition 1: unimodality under concave accuracy}\label{app:prop:unimodal}

\paragraph{Proposition 1.} Assume (A1), (A2), and $A''(k)\le0$ on $(0,K)$; write
$J(k)=a+bk$ and define the stationarity function
\begin{equation}\label{app:eq:N}
  N(k)=bA(k)-J(k)A'(k)=A(k)^{2}E'(k)
  \qquad [\mathrm{J\,query^{-1}\,passage^{-1}}].
\end{equation}
Then:
\begin{enumerate}
  \item[(i)] $N'(k)=-(a+bk)A''(k)\ge0$, so $N$ is nondecreasing on $[0,K]$, and
  strictly increasing if $A''<0$ on $(0,K)$.
  \item[(ii)] $E'=N/A^{2}$ changes sign at most once on $(0,K)$, and any sign
  change is from negative to positive. Hence $E$ is nonincreasing then
  nondecreasing: it is unimodal (quasiconvex) and has no interior local maximum.
  If $A''<0$, $E$ has at most one interior stationary point, and when it exists it
  is the unique global minimiser of $E$ on $[0,K]$.
  \item[(iii)] If $A''<0$, an interior minimiser exists if and only if
  \begin{equation}\label{app:eq:interior}
    bA(0)<aA'(0)
    \quad\text{and}\quad
    bA(K)>J(K)A'(K);
  \end{equation}
  otherwise the minimiser is $k=0$ (if the first inequality fails) or $k=K$ (if
  the second fails).
\end{enumerate}

\paragraph{Proof.} (i) Differentiate \eqref{app:eq:N} using (A1):
\begin{equation}\label{app:eq:Nprime}
  N'(k)=bA'(k)-\left[bA'(k)+(a+bk)A''(k)\right]=-(a+bk)A''(k),
\end{equation}
and $a+bk>0$ by (A1), so $N'\ge0$; if $A''<0$ then $N'>0$.

(ii) By \eqref{app:eq:N}, $E'=N/A^{2}$ and $A^{2}>0$, so $E'$ has the sign of
$N$. A nondecreasing function crosses zero at most once, and when it does the
crossing is from negative to positive; hence there is a point $\bar k$ (possibly
$0$ or $K$, possibly an interval when $N$ is flat at zero) with $E'\le0$ for
$k<\bar k$ and $E'\ge0$ for $k>\bar k$. That is unimodality, and no interior local
maximum is possible. If $A''<0$ then $N$ is strictly increasing, so its zero set
is empty or a single point $k^{*}$. At such a point,
\begin{equation}\label{app:eq:Esecond}
  E''(k^{*})=\frac{N'(k^{*})A(k^{*})^{2}-2A(k^{*})A'(k^{*})N(k^{*})}{A(k^{*})^{4}}
  =\frac{N'(k^{*})}{A(k^{*})^{2}}>0,
\end{equation}
because $N(k^{*})=0$; the point is a strict local minimum, and by the sign pattern
it is the unique global minimiser on $[0,K]$.

(iii) If $A''<0$ then $N$ is continuous and strictly increasing. An interior zero
exists if and only if $N(0)<0<N(K)$; substituting \eqref{app:eq:N} gives
$bA(0)-aA'(0)<0$ and $bA(K)-(a+bK)A'(K)>0$, which is \eqref{app:eq:interior}.
If $N(0)\ge0$ then $N>0$ on $(0,K]$ and $E$ is nondecreasing, so the minimiser is
$k=0$; if $N(K)\le0$ then $N<0$ on $[0,K)$ and $E$ is nonincreasing, so the
minimiser is $k=K$. \hfill$\square$

\paragraph{Two remarks.} First, the proof never uses $A'\ge0$: monotonicity of
$N$ depends only on the sign of $A''$. The proposition therefore also covers a
concave decline after an accuracy peak, which is the qualitative shape of the two
non-monotone P4 arms (a peak followed by a decline), although concavity of the
decline is not tested there. Second, the hypothesis is testable on the grid: the discrete
analogue of $A''\le0$ is that the secant slopes between adjacent grid depths are
non-increasing, which is what the companion implementation checks. The hypothesis
is not a description of every accuracy curve; where it fails, Proposition~1 gives
no global statement and the identity remains local.

\paragraph{Existence in elasticity form.} Conditions \eqref{app:eq:interior} can
be read as bounds on the elasticities. The first is
$b/a<A'(0)/A(0)$, i.e.\ the relative accuracy gain per passage at zero must exceed
the relative energy cost per passage. The second is
$\varepsilon_{A}(K)<\varepsilon_{J}(K)=bK/(a+bK)<1$, i.e.\ the accuracy
elasticity at the top of the grid must be below the energy elasticity. Both are
properties of the fitted curves, not of the decision rule, and both are evaluated
without the accuracy gate of (A3).

\subsection{Corollary 1: comparative statics of the optimum}\label{app:cor:statics}

\paragraph{Corollary 1.} Assume (A1), (A2), $A'(k^{*})>0$ and $A''(k^{*})<0$ at an
interior stationary point $k^{*}=k^{*}(a,b)$ of $E$. Then $k^{*}$ is continuously
differentiable in $(a,b)$ in a neighbourhood, and
\begin{equation}\label{app:eq:statics}
  \frac{\partial k^{*}}{\partial a}
  =\frac{-A'(k^{*})}{(a+bk^{*})A''(k^{*})}>0,
  \qquad
  \frac{\partial k^{*}}{\partial b}
  =\frac{A(k^{*})-k^{*}A'(k^{*})}{(a+bk^{*})A''(k^{*})}
  =-\frac{aA(k^{*})}{(a+bk^{*})^{2}\left(-A''(k^{*})\right)}<0 .
\end{equation}

\paragraph{Proof.} By Lemma~1 the stationary point solves
$F(k,a,b)=bA(k)-(a+bk)A'(k)=0$, and $F$ is continuously differentiable under
(A2). Its partials are
\begin{equation}\label{app:eq:Fpartials}
  F_{k}=-(a+bk)A''(k),
  \qquad
  F_{a}=-A'(k),
  \qquad
  F_{b}=A(k)-kA'(k).
\end{equation}
At $k^{*}$, $F_{k}=-(a+bk^{*})A''(k^{*})>0$ because $A''(k^{*})<0$; this is also
the second-order condition \eqref{app:eq:Esecond}. The implicit function theorem
therefore applies, and
\begin{equation}\label{app:eq:IFT}
  \frac{\partial k^{*}}{\partial a}=-\frac{F_{a}}{F_{k}},
  \qquad
  \frac{\partial k^{*}}{\partial b}=-\frac{F_{b}}{F_{k}} .
\end{equation}
Substituting \eqref{app:eq:Fpartials} gives the first equality in each line of
\eqref{app:eq:statics}. The first sign is immediate since $A'(k^{*})>0$ and the
denominator is negative. For the second, apply the identity
\eqref{app:eq:identity-iff} at $k^{*}$:
$k^{*}A'(k^{*})/A(k^{*})=\varepsilon_{J}(k^{*})=bk^{*}/(a+bk^{*})<1$, hence
\begin{equation}\label{app:eq:Fbpositive}
  A(k^{*})-k^{*}A'(k^{*})
  =A(k^{*})\left(1-\varepsilon_{J}(k^{*})\right)
  =\frac{aA(k^{*})}{a+bk^{*}}>0,
\end{equation}
and the denominator of the second line is negative, so
$\partial k^{*}/\partial b<0$. \hfill$\square$

\paragraph{Reading.} A larger fixed cost per query makes depth relatively cheap
and pushes the unconstrained optimum deeper; a larger per-passage cost makes
depth relatively expensive and pushes it shallower. The corollary is local and
applies only when the accuracy gate is slack: when the gate binds, the policy
depth is the smallest feasible grid depth and is locally insensitive to $(a,b)$,
so the derivative is zero in the one-sided sense at that corner. It is also
silent about the magnitude, because $A'(k^{*})$ and $A''(k^{*})$ are estimated
from the same finite sample as $A(k^{*})$; the sign is a prediction, the size is
not.

\subsection{Corollary 2: when the optimum is the boundary}\label{app:cor:boundary}

\paragraph{Corollary 2.} Assume (A1) and the stylized coverage family
\begin{equation}\label{app:eq:coverage}
  A(k)=1-(1-a_{0})(1-p)^{k},
  \qquad
  a_{0}=A(0)\in(0,1),
  \qquad
  p\in(0,1),
\end{equation}
so that $A(k)\in(a_{0},1)$ and (A2) holds on $[0,\infty)$. Let
$L=-\ln(1-p)>0$. Then $A'(k)=L\left(1-A(k)\right)>0$ and
$A''(k)=-L^{2}\left(1-A(k)\right)<0$, Proposition~\ref{app:prop:unimodal}
applies, and
\begin{equation}\label{app:eq:boundary}
  k^{*}=0
  \iff
  \frac{b}{a}\ge\frac{(1-a_{0})\left(-\ln(1-p)\right)}{a_{0}}
  =\frac{1-a_{0}}{a_{0}}\,L .
\end{equation}
When the inequality is reversed, the unique interior optimum solves
\begin{equation}\label{app:eq:coverage-optimum}
  (1-a_{0})(1-p)^{k^{*}}=\frac{b}{aL+b+bLk^{*}} .
\end{equation}

\paragraph{Proof.} Differentiating \eqref{app:eq:coverage} gives
$A'(k)=-(1-a_{0})(1-p)^{k}\ln(1-p)=L\left(1-A(k)\right)$ and
$A''(k)=-L^{2}\left(1-A(k)\right)<0$, so $A$ is strictly increasing and strictly
concave and (A2) holds. Evaluate \eqref{app:eq:N} at zero:
\begin{equation}\label{app:eq:N0}
  N(0)=bA(0)-aA'(0)=ba_{0}-a(1-a_{0})L
  =a_{0}\left[b-a\frac{1-a_{0}}{a_{0}}L\right].
\end{equation}
By Proposition~\ref{app:prop:unimodal}(i), $N$ is strictly increasing. As
$k\to\infty$, $A(k)\to1$ and $(a+bk)(1-A(k))\to0$ because the exponential decay
dominates the linear factor, so
\begin{equation}\label{app:eq:Nlimit}
  N(k)=bA(k)-(a+bk)L\left(1-A(k)\right)\longrightarrow b>0 .
\end{equation}
If $N(0)\ge0$, then $N>0$ on $(0,\infty)$ and $E$ is strictly increasing, so
$k^{*}=0$; by \eqref{app:eq:N0} this is exactly the inequality in
\eqref{app:eq:boundary}. If $N(0)<0$, then \eqref{app:eq:Nlimit} and strict
monotonicity give a unique zero $k^{*}\in(0,\infty)$, which is the unique
minimiser by Proposition~\ref{app:prop:unimodal}(ii). Setting $N(k^{*})=0$ and
writing $s=1-A(k^{*})=(1-a_{0})(1-p)^{k^{*}}$ gives
$b(1-s)=(a+bk^{*})Ls$, i.e.\ $b=s\left(aL+b+bLk^{*}\right)$, which is
\eqref{app:eq:coverage-optimum}. \hfill$\square$

\paragraph{Scope.} The corollary is a statement about one stylized family, not
about the observed curves. The P4 validation rejects the family as a general
description of the four arms ($R^{2}$ between $0.097$ and $0.879$), so its role
is to show that a boundary optimum is a possible outcome of diminishing returns,
and to give the threshold in closed form. It does not certify any observed
$k^{*}$, and it is not used to choose a depth.

\subsection{Proposition 2: the price of static depth}\label{app:prop:price}

\paragraph{Proposition 2.} Assume (A1). Let $\pi$ be any per-question depth
assignment with values in $[0,K]$, with population energy
$J_{\pi}=\mathbb{E}_{q}[J(\pi(q))]$ and accuracy
$A_{\pi}=\mathbb{E}_{q}[A(\pi(q))]$, and let $k_{s}$ be a static depth with
accuracy $A_{s}=A(k_{s})$. If $A_{\pi}\ge A_{s}$, then
\begin{equation}\label{app:eq:ceiling}
  \frac{E_{\pi}}{E_{s}}
  =\frac{J_{\pi}}{J(k_{s})}\cdot\frac{A_{s}}{A_{\pi}}
  \le\frac{a+bK}{a}
  =\frac{1}{1-\varepsilon_{J}(K)} ,
\end{equation}
and the largest saving any such assignment can deliver is the variable energy
share at the reference depth,
\begin{equation}\label{app:eq:ceiling-saving}
  1-\frac{E_{\pi}}{E_{s}}\le\frac{bK}{a+bK}=\varepsilon_{J}(K)<1 .
\end{equation}
The bound is sharp. For $\varepsilon\in(0,1)$ let a fraction $1-\varepsilon$ of
questions be easy, answered correctly at every depth including $k=0$, and a
fraction $\varepsilon$ be hard, answered correctly only at depths $k\ge K$. Then
the best static depth at accuracy $1$ is $k_{s}=K$, the hindsight oracle
$\pi_{\mathrm{or}}$ assigns easy questions to $0$ and hard questions to $K$, and
\begin{equation}\label{app:eq:rho}
  \rho(\varepsilon)
  =\frac{E_{\mathrm{static}}}{E_{\mathrm{oracle}}}
  =\frac{a+bK}{a+\varepsilon bK}
  \;\longrightarrow\;
  1+\frac{bK}{a}
  =\frac{1}{1-\varepsilon_{J}(K)}
  \qquad(\varepsilon\to0).
\end{equation}
The ratio is written here as static over oracle, so that the quantity is the
\emph{price} of fixing one depth; written in the opposite order (oracle over
static) the same limit is its reciprocal $a/(a+bK)<1$.

\paragraph{Proof.} By (A1), $J(\pi(q))\in[a,a+bK]$ for every question, so
$a\le J_{\pi}\le a+bK$ and $J(k_{s})\ge a$; hence
$J_{\pi}/J(k_{s})\le(a+bK)/a$. Since $A_{\pi}\ge A_{s}>0$, the factor
$A_{s}/A_{\pi}$ is at most one, which proves \eqref{app:eq:ceiling}. For
\eqref{app:eq:ceiling-saving}, subtract the reciprocal:
$1-a/(a+bK)=bK/(a+bK)=\varepsilon_{J}(K)$ by (A1).

For sharpness, let $q$ be easy with probability $1-\varepsilon$ and hard with
probability $\varepsilon$. Every static depth $k<K$ answers only the easy
questions, so its accuracy is $1-\varepsilon<1$; the only static depths with
accuracy $1$ are $k\ge K$, and among them $J$ is minimised at $k=K$. Hence
$E_{\mathrm{static}}=a+bK$. The oracle answers every question correctly, so its
accuracy is $1$ and its energy is
$(1-\varepsilon)a+\varepsilon(a+bK)=a+\varepsilon bK$, giving
$E_{\mathrm{oracle}}=a+\varepsilon bK$ and the first equality in
\eqref{app:eq:rho}. The limit is immediate, and the last equality repeats
$bK/(a+bK)=\varepsilon_{J}(K)$. \hfill$\square$

\paragraph{Scope.} Proposition~\ref{app:prop:price} is an upper bound with a
sharpness example, not a lower bound on the value of adaptivity and not a claim
about the observed workload. Two consequences matter for the paper. First, the
ceiling is a property of the affine energy model alone: it caps \emph{any}
allocation, learned or oracle, at matched accuracy, and it is the formal reason
the oracle gain is reported as a fraction of the static energy rather than as an
unbounded promise. Second, the example shows that the ceiling is approached when
a small hard stratum forces the whole workload to the deepest setting; when the
stratum is large ($\varepsilon\to1$), the ratio tends to one and adaptivity is
worth nothing. The empirical matched-oracle gain reported for P3 is far below the
P1 ceiling of $\varepsilon_{J}(32)=0.7517$, as it must be, and the pilot oracle
figure of $78.8\%$ is below the pilot ceiling
$bK/(a+bK)=0.8227$ computed from $a=15.74$ and $b=2.283$; both are consistency
checks on registered quantities, not new evidence.

\subsection{Proposition 3: the cost of certification}\label{app:prop:cert}

Certification has two separate costs: the energy of the experiment (part a) and
the number of questions needed for the statistical statement (part b). They are
stated separately because the first is a deterministic break-even and the second
is a sampling requirement, and only their product is the cost of a certified
depth.

\paragraph{Proposition 3(a): economic break-even.} Assume (A1), a deployment of
$N$ queries drawn from the same distribution as the experiment, one depth used
for all $N$ queries, no discounting, and a one-off certification cost
$C_{\mathrm{cert}}>0$ in joules. Let $k_{0}=K$ be the reference depth and
$k^{*}$ the candidate. Define the per-query break-even and the
per-correct-answer break-even by
\begin{equation}\label{app:eq:nstar}
  N^{*}_{J}=\frac{C_{\mathrm{cert}}}{J(k_{0})-J(k^{*})},
  \qquad
  N^{*}_{A}=\frac{C_{\mathrm{cert}}}{A(k^{*})\left[E(k_{0})-E(k^{*})\right]},
\end{equation}
both in queries and both defined only when the denominator is positive. Then
\begin{equation}\label{app:eq:nstar-identity}
  A(k^{*})\left[E(k_{0})-E(k^{*})\right]
  =\frac{A(k^{*})}{A(k_{0})}J(k_{0})-J(k^{*}),
\end{equation}
so $N^{*}_{A}=N^{*}_{J}$ if and only if $A(k^{*})=A(k_{0})$, and otherwise
\begin{equation}\label{app:eq:beta}
  \beta:=\frac{N^{*}_{A}}{N^{*}_{J}}
  =\frac{J(k_{0})-J(k^{*})}{\dfrac{A(k^{*})}{A(k_{0})}J(k_{0})-J(k^{*})},
\end{equation}
with $\beta\ge1$ if and only if $A(k^{*})\le A(k_{0})$.

\paragraph{Proof.} Equation \eqref{app:eq:nstar-identity} is
$A(k^{*})E(k_{0})-A(k^{*})E(k^{*})=A(k^{*})J(k_{0})/A(k_{0})-J(k^{*})$, using
$A(k^{*})E(k^{*})=J(k^{*})$. Dividing the two lines of \eqref{app:eq:nstar}
gives \eqref{app:eq:beta}. For the sign, compare denominators:
\begin{equation}\label{app:eq:beta-sign}
  \left[J(k_{0})-J(k^{*})\right]
  -\left[\frac{A(k^{*})}{A(k_{0})}J(k_{0})-J(k^{*})\right]
  =J(k_{0})\left(1-\frac{A(k^{*})}{A(k_{0})}\right),
\end{equation}
which is non-negative exactly when $A(k^{*})\le A(k_{0})$. Both denominators are
positive by hypothesis, so the inequality passes to $\beta$.
\hfill$\square$

\paragraph{Interpretation and the approximation.} $N^{*}_{J}$ is the number of
deployment queries at which the energy saved by switching depth equals the
certificate: after $N$ queries at $k^{*}$ instead of $k_{0}$ the saving is
$N\left[J(k_{0})-J(k^{*})\right]$, and it repays $C_{\mathrm{cert}}$ at
$N=N^{*}_{J}$. $N^{*}_{A}$ answers a different question: it values the candidate's
correct-answer output at the per-correct-answer saving $E(k_{0})-E(k^{*})$, so $N$
candidate-depth queries produce $N A(k^{*})$ correct answers and a cumulative
saving $N A(k^{*})\left[E(k_{0})-E(k^{*})\right]$, which equals
$C_{\mathrm{cert}}$ at $N=N^{*}_{A}$. The two coincide only at
equal accuracy. When $A(k^{*})<A(k_{0})$, the paper's per-correct-answer formula
is therefore an approximation to the deployment-queries reading, and it errs on
the conservative side ($\beta>1$): it reports a longer break-even because it also
requires the candidate to reproduce the correct-answer output that its lower
accuracy would otherwise lose. The direction is known but the formula is not
exact, which is why both numbers are reported rather than one.

\paragraph{P1 diagnostic.} With the registered P1 values
$C_{\mathrm{cert}}=311{,}992.064$ J, $J(4)=36.2537$ J/query,
$J(32)=94.9450$ J/query, $A(4)=0.450$, and $A(32)=0.493$, one obtains
$N^{*}_{J}=5{,}315.8$ queries, $N^{*}_{A}=6{,}189.1$ queries, and
$\beta=1.1643$. The certification cost is itself a sum over the grid,
$C_{\mathrm{cert}}=n\sum_{k\in\mathcal{K}}J(k)=1{,}000\times311.992$ J, so the
sweep energy and the reported break-even use the same accounting.

\paragraph{Proposition 3(b): statistical sample complexity.} Assume (A5), let
$m$ be the accuracy margin (written $\delta$ in the main text, and renamed here
to avoid a clash with the error probability), and let
$\Delta=A(k_{0})-A(k^{*})>0$. For the paired differences
$D_{q}=C_{qk_{0}}-C_{qk^{*}}\in\{-1,0,1\}$ with
$\mathbb{E}[D_{q}]=\Delta$, put
$t_{n}=\sqrt{2\ln(1/\alpha)/n}$ and declare non-inferiority of $k^{*}$ when
\begin{equation}\label{app:eq:declare}
  \bar D+t_{n}<m .
\end{equation}
Then:
\begin{enumerate}
  \item[(i)] \eqref{app:eq:declare} is a level-$\alpha$ test of
  $H_{0}:\Delta\ge m$.
  \item[(ii)] At a true difference $\Delta=m_{1}<m$ the probability of failing to
  declare non-inferiority is at most
  \begin{equation}\label{app:eq:hoeffding}
    \Pr\left(\text{fail to declare}\mid \Delta=m_{1}\right)
    \le\exp\left(-\frac{n\left(\left(m-m_{1}\right)-t_{n}\right)^{2}}{2}\right),
  \end{equation}
  so $n\ge8\max\left(\ln(1/\alpha),\ln(1/\delta_{\mathrm{err}})\right)/(m-m_{1})^{2}$
  paired questions suffice for a failure probability at most
  $\delta_{\mathrm{err}}$.
  \item[(iii)] Selecting the more accurate of two depths at a true gap $\Delta$
  needs $n\ge2\ln(1/\delta_{\mathrm{err}})/\Delta^{2}$, and selecting among $d$
  depths whose smallest gap to the optimum is $\Delta_{\min}$ needs
  $n\ge2\ln\left((d-1)/\delta_{\mathrm{err}}\right)/\Delta_{\min}^{2}$ by a union
  bound; in both cases the correct depth is selected with probability at least
  $1-\delta_{\mathrm{err}}$.
\end{enumerate}

\paragraph{Proof sketch.} $D_{q}$ is bounded in $[-1,1]$ and has range $2$, so
Hoeffding's inequality gives, for every $u>0$,
\begin{equation}\label{app:eq:two-tails}
  \Pr\left(\bar D-\Delta\le-u\right)\le\exp\left(-\frac{nu^{2}}{2}\right),
  \qquad
  \Pr\left(\bar D-\Delta\ge u\right)\le\exp\left(-\frac{nu^{2}}{2}\right).
\end{equation}
(i) Under $H_{0}$, $\Delta\ge m$, so declaring non-inferiority implies
$\bar D-\Delta<-t_{n}$, which has probability at most
$\exp(-nt_{n}^{2}/2)=\alpha$.

(ii) Failure to declare means $\bar D+t_{n}\ge m$, hence
$\bar D-\Delta\ge(m-m_{1})-t_{n}$, and the upper tail of
\eqref{app:eq:two-tails} applies. If $t_{n}\le(m-m_{1})/2$, which holds once
$n\ge8\ln(1/\alpha)/(m-m_{1})^{2}$, the exponent is at most
$-n(m-m_{1})^{2}/8$, and the stated $n$ makes it at most
$\ln(1/\delta_{\mathrm{err}})$.

(iii) The selection error is
$\Pr(\bar D\le0)=\Pr(\bar D-\Delta\le-\Delta)\le\exp(-n\Delta^{2}/2)$, and the
union bound over the $d-1$ alternatives gives the second statement.
\hfill$\square$

\paragraph{Scope and consequences.} The bound is distribution-free and makes no
normal approximation, but it is deliberately conservative: it ignores the
pairing, i.e.\ the discordance rate, which is what a McNemar calculation uses.
The tighter normal-approximation requirement scales like
$\psi/(m-m_{1})^{2}$ at discordance $\psi$, so at a fixed difference the observed
P1 discordance of $13.1\%$ requires roughly $1.6$ times the pairs that the
planned $8\%$ implied, and the registered frontier puts the 5 pp requirement at
$n=7{,}138$ pairs. Three consequences are used in the main text. First, the
sample complexity diverges as $m_{1}\uparrow m$: a margin that the data show is
several percentage points wide cannot be resolved by a design powered at a
near-zero difference, which is the formal content of the P1 precision failure.
Second, the energy side does not drive the requirement: per-query energy has a P1
repeat CV of $2.81\%$ at $k=8$, so the relative standard error of a mean energy
at $n=1{,}000$ is about $2.81\%/\sqrt{1{,}000}<0.1\%$, while the accuracy
contrast at 5 pp remains unresolved at the same $n$. The accuracy side is the
binding one. Third, the certification cost is linear in the statistical
requirement, because a sweep of $n$ questions costs
$C_{\mathrm{cert}}(n)=n\,c_{\mathrm{sweep}}$ with
$c_{\mathrm{sweep}}=\sum_{k\in\mathcal{K}}J(k)$ in J/query, so
\begin{equation}\label{app:eq:nstar-n}
  N^{*}_{J}(n)=\frac{n\,c_{\mathrm{sweep}}}{J(k_{0})-J(k^{*})},
\end{equation}
which grows linearly in $n$. The statistical requirement and the economic
break-even therefore cannot be chosen independently, and a depth that needs a
large experiment needs a proportionally larger deployment to justify it.

\subsection{Grid resolution and the H3 diagnostic}\label{app:grid}

The identity of Lemma~1 is a statement about a derivative. On the grid two
approximations intervene, and separating them is what makes the H3 result
readable.

\emph{First, the derivative is a secant.} The reported accuracy elasticity at a
grid depth is a finite difference over the neighbouring depths, so the identity
at a grid point is the statement that the secant elasticity equals the affine
elasticity. At $k=4$ the central difference uses $k=2$ and $k=8$, a step of three
passages; at $k=8$ it uses $k=4$ and $k=16$, a step of six; at the edges the
difference is one-sided, with step one at $k=1$ and step sixteen at $k=32$. A
central difference approximates $A'(k)$ to second order in the step and a
one-sided difference to first order, so a grid gap is a discretisation of
$\Delta(k)$, not a different quantity.

\emph{Second, the decision object is constrained.} The policy depth of (A3) is the
energy minimiser inside the feasible set $\mathcal{F}_{m}$, whereas $k^{*}$ in
Lemma~1 is the unconstrained stationary point of $E$. The two objects coincide
only when the accuracy gate is slack. In P1 the gate binds: the 5 pp rule selects
$k=4$ because it is the smallest feasible depth, not because $E$ is stationary
there. Testing the identity at the gate-selected depth is therefore a test of
local stationarity at a point that the theory does not expect to be stationary,
and the main text reports it as such.

\emph{The P1 numbers.} With $a=24.0759$ J/query and
$b=2.2772$ J per passage per query, the affine elasticity is
$\varepsilon_{J}(4)=0.2745$, the secant elasticity is
$\varepsilon_{A}(4)=0.0741$, and the gap is $\Delta(4)=+0.2004$ with a 95\%
interval of $[0.1677,0.2320]$ that excludes zero. The positive sign says that
$E$ is increasing at $k=4$: the last passages bought at that depth cost more
energy per correct answer than they return. The monotone-cubic interpolant of the
same seven means has its unconstrained stationary point at $k^{*}=1.2144$ with
$E=67.1408$ J/correct and a gap of $-1.9\times10^{-7}$, which is zero to numerical
precision. Since $1.2144<4$, a positive gap at $k=4$ is exactly what the theory
requires. The H3 test asked whether the gate-selected depth is stationary, and
the answer is no; that is a statement about the selected depth, not a failure of
the identity.

\emph{What the interpolation does and does not establish.} The near-zero gap at
$k^{*}=1.2144$ is close to a consistency check on the solver: any smooth
interpolant, minimised numerically, will report $E'\approx0$ at its own
minimiser, so the zero is not independent evidence for the location. The
independent content is the sign at $k=4$, which is estimated from the data and
has an interval that excludes zero, and the flatness of $E$ between $k=1$ and
$k=2$: the two values are $73.71$ and $73.33$ J/correct, a difference of $0.5\%$,
with overlapping percentile intervals $[68.2,79.9]$ and $[68.1,79.3]$. The
existence of an interior unconstrained optimum is therefore supported on this
arm, while its location inside $[1,2]$ is not identified. The defensible
statement is that $k=4$ is over-retrieval relative to an unconstrained optimum
near $k=1.2$, not that $1.214$ is the optimal depth. Depth $0$ is excluded from
the identity because $\varepsilon$ is degenerate there and because $A(0)=0.001$
rests on a single correct answer.

\subsection{What is not claimed}\label{app:notclaimed}

The appendix is deliberately narrower than a general theory of retrieval depth.
It does not claim the following.

\begin{itemize}
  \item \emph{A universal $k^{*}$.} Every number quoted here is a property of one
  reader, one retriever, one corpus, one grid, and one hardware condition. The
  comparative statics are signs, not magnitudes, and they are derived for the
  unconstrained ratio.
  \item \emph{That accuracy is concave in general.} Proposition~\ref{app:prop:unimodal}
  is conditional. The P4 validation rejects the coverage family on three of four
  arms and finds non-monotone curves on two, so the hypothesis is tested where it
  is used and is not a maintained assumption about RAG.
  \item \emph{That a zero gap certifies optimality.} The identity is necessary,
  not sufficient, as the convex counterexample in Section~\ref{app:lemma1} shows.
  A gap interval that contains zero is consistent with stationarity; it does not
  establish a minimum, and it cannot speak for a boundary depth.
  \item \emph{That the grid argmin is the continuous argmin.} The grid is coarse
  and uneven, the interpolant is a modelling choice, and the policy object is the
  constrained minimiser of (A3). The two optima are reported separately and are
  never substituted for one another.
  \item \emph{That the ceiling of Proposition~\ref{app:prop:price} is attained by
  a deployable policy.} It is a bound with a sharpness example in a stylized
  mixture. A learned policy is a separate object that must be trained, selected,
  and evaluated on disjoint splits, and the P3 result shows that a realizable
  policy can capture none of the gap.
  \item \emph{That $N^{*}$ is a deployment guarantee.} It is a deterministic
  break-even: no discounting, no uncertainty in the saving, no cost of capital,
  and a certificate that is assumed to settle the question once paid.
  \item \emph{Any claim outside the measured condition.} No carbon claim, no
  claim about other hardware, readers, or retrievers, and no claim that the
  exploratory arms enter the theory.
  \item \emph{Any re-adjudication.} The appendix proves statements; it does not
  change the status of H3 or of any gate. H3 remains adjudicated at $k=4$ and
  remains unsupported.
\end{itemize}

%
%  The manuscript preamble loads amsmath, amssymb, and booktabs but not amsthm,
%  so results are set as \paragraph run-in headings and every proof ends with
%  \hfill$\square$. No preamble change is required, and this file can be built
%  on its own inside the same class.
%
%  Scope: this appendix proves only the statements used in the main text.
%  Values marked "P1 diagnostic" are registered outputs of
%  analysis/squad_qwen4b_p1/ and analysis/identity_continuous.json; they are
%  quoted, not re-estimated, and no gate is re-adjudicated here.
% ===========================================================================
\section{Formal theory of the depth decision}\label{app:theory}

This appendix states and proves the six results the manuscript relies on: the
first-order identity (Lemma~1), the shape condition that makes the identity
useful (Proposition~1), the comparative statics of the optimum (Corollary~1),
the boundary condition for a zero optimum (Corollary~2), the ceiling on the value
of per-question allocation (Proposition~2), and the cost of certifying a depth
(Proposition~3). Section~\ref{app:grid} separates the grid identity from the
continuous optimum, and Section~\ref{app:notclaimed} lists what the theory does
not license. Every result is conditional on the assumptions in
Section~\ref{app:assumptions}; none of them is a statement about the reader, the
retriever, or the hardware as such.

\subsection{Setting and assumptions}\label{app:assumptions}

Depth $k$ is a count of passages and ranges over the closed interval $[0,K]$ with
$K=32$. The measured grid is
\begin{equation}\label{app:eq:grid}
  \mathcal{K}=\{0,1,2,4,8,16,32\},
\end{equation}
so the grid is coarse and unevenly spaced, with gaps of $1$, $1$, $2$, $4$, $8$,
and $16$ passages. Mean energy per query is $J(k)$ in joules per query, accuracy
is $A(k)$ as a proportion, and energy per correct answer is the ratio estimand of
the main text,
\begin{equation}\label{app:eq:ratio}
  E(k)=\frac{J(k)}{A(k)}
  \qquad [\mathrm{J/correct}].
\end{equation}
The two elasticities and their gap are used throughout:
\begin{equation}\label{app:eq:elasticities}
  \varepsilon_{J}(k)=\frac{kJ'(k)}{J(k)},
  \qquad
  \varepsilon_{A}(k)=\frac{kA'(k)}{A(k)},
  \qquad
  \Delta(k)=\varepsilon_{J}(k)-\varepsilon_{A}(k),
\end{equation}
both dimensionless for $k>0$. The assumptions are as follows.

\begin{enumerate}
  \item[(A1)] \emph{Affine energy.} $J(k)=a+bk$ for $k\in[0,K]$, with $a>0$ in
  J/query and $b>0$ in J per passage per query. Equivalently
  $\varepsilon_{J}(k)=bk/(a+bk)\in(0,1)$ and $\varepsilon_{J}$ is strictly
  increasing in $k$.

  \item[(A2)] \emph{Positive accuracy.} $A(k)\in(0,1]$ for $k\in[0,K]$, and $A$ is
  twice differentiable wherever a derivative appears. Curvature signs are stated
  per result and are never assumed globally without being named.

  \item[(A3)] \emph{Decision object.} The policy quantity is the constrained
  minimiser over the grid, $k^{*}_{m}=\arg\min_{k\in\mathcal{F}_{m}}E(k)$ with
  $\mathcal{F}_{m}=\{k\in\mathcal{K}:A(k)\ge A(K)-m\}$, which is the object of the
  main text. Statements about a continuous $k^{*}$ are statements about the
  extension in (A4), not about unmeasured depths.

  \item[(A4)] \emph{Continuum extension.} Off the grid, $A$ denotes a twice
  differentiable interpolant of the seven grid means (the registered analysis
  uses a monotone cubic). The interpolant is a modelling device: it makes
  $E'(k)$ well defined, and every off-grid number inherits its smoothness
  assumption.

  \item[(A5)] \emph{Sampling.} Questions are independent draws from a fixed
  distribution, each question is answered at every depth in $\mathcal{K}$, and
  correctness is binary, $C_{qk}\in\{0,1\}$ with $\mathbb{E}[C_{qk}]=A(k)$. This
  assumption is used only in Proposition~\ref{app:prop:cert}(b).

  \item[(A6)] \emph{Calibration.} Calibrated energy is a positive depth-common
  multiple of the measured energy, $J_{\mathrm{cal}}(k)=fJ(k)$ with $f>0$. The
  calibration rule of the main text is all-or-nothing across depths, so this is
  the regime in which absolute energy is reported.
\end{enumerate}

Assumption (A6) has an immediate consequence that justifies using calibrated
energy in the theory: $E_{\mathrm{cal}}=fE$, so the minimiser, every ratio of two
depths, and every elasticity $\varepsilon_{J}$ are unchanged, because $f$ cancels
in $kJ'_{\mathrm{cal}}/J_{\mathrm{cal}}$. A depth-dependent calibration error
would not cancel, which is why the measurement protocol gates on a per-depth
factor and applies it only if all depths pass.

\begin{table}[pos=htbp]
  \centering
  \caption{Summary of the appendix. Each result is conditional on the
  hypotheses named in its statement; none of them is a claim about a particular
  reader, corpus, or GPU.}
  \label{tab:theorymap}
  \footnotesize
  \setlength{\tabcolsep}{4pt}
  \begin{tabular}{@{}>{\raggedright\arraybackslash}p{1.9cm}>{\raggedright\arraybackslash}p{7.4cm}>{\raggedright\arraybackslash}p{5.0cm}@{}}
    \toprule
    \textbf{Result} & \textbf{What is proved} & \textbf{What it licenses} \\
    \midrule
    Lemma~1 & $E'(k)=0$ iff the two elasticities agree, for $k>0$ &
    A local diagnostic; necessary, not sufficient \\
    Proposition~1 & Concave $A$ makes the stationarity function monotone, so $E$
    is unimodal & Existence and uniqueness of an interior optimum, and a test of
    the hypothesis \\
    Corollary~1 & $\partial k^{*}/\partial a>0$ and $\partial k^{*}/\partial b<0$
    at a slack-gate interior optimum & Directional, qualitative predictions only \\
    Corollary~2 & A closed-form $b/a$ threshold for $k^{*}=0$ under coverage &
    Boundary optima are possible; a check on the coverage family only \\
    Proposition~2 & Any allocation at matched accuracy satisfies
    $E_{\pi}/E_{s}\le 1+bK/a$ & A ceiling on the value of adaptivity, with a
    sharpness example \\
    Proposition~3 & Break-even volume and sample complexity of certification &
    A cost, not a guarantee; the accuracy side dominates \\
    \bottomrule
  \end{tabular}
\end{table}

\subsection{Lemma 1: the first-order identity}\label{app:lemma1}

\paragraph{Lemma 1.} Under (A1) and (A2), for every $k\in(0,K]$,
\begin{equation}\label{app:eq:identity}
  E'(k)=\frac{bA(k)-J(k)A'(k)}{A(k)^{2}}
  =\frac{J(k)}{kA(k)}\,\Delta(k),
\end{equation}
and therefore
\begin{equation}\label{app:eq:identity-iff}
  E'(k)=0
  \iff
  bA(k)=J(k)A'(k)
  \iff
  \frac{kJ'(k)}{J(k)}=\frac{kA'(k)}{A(k)}
  \iff
  \Delta(k)=0 .
\end{equation}
At $k=0$ only the unmultiplied form is meaningful,
$E'(0)=[bA(0)-aA'(0)]/A(0)^{2}$, because the elasticity form is identically
$0=0$ there.

\paragraph{Proof.} By the quotient rule and (A1),
\begin{equation}\label{app:eq:proof-quotient}
  E'(k)=\frac{J'(k)A(k)-J(k)A'(k)}{A(k)^{2}}
  =\frac{bA(k)-J(k)A'(k)}{A(k)^{2}},
\end{equation}
which is the first equality in \eqref{app:eq:identity}; $A(k)>0$ by (A2), so the
expression is defined. For the second equality, write
$E=J/A$ and differentiate logarithmically, $E'/E=J'/J-A'/A$. Multiplying by $k$
gives $kE'(k)/E(k)=\Delta(k)$ and hence
\begin{equation}\label{app:eq:proof-second}
  E'(k)=\frac{E(k)}{k}\Delta(k)=\frac{J(k)}{kA(k)}\Delta(k),
\end{equation}
which agrees with \eqref{app:eq:proof-quotient} after multiplying both sides by
$kA(k)^{2}$. Since $J(k)>0$, $A(k)>0$, and $k>0$, the prefactor $J/(kA)$ in
\eqref{app:eq:proof-second} is strictly positive, so $E'(k)=0$ holds if and only
if $\Delta(k)=0$; the middle equivalence in \eqref{app:eq:identity-iff} is the
same statement multiplied by $k>0$ and divided by the positive quantities $J$
and $A$. \hfill$\square$

\paragraph{Necessity, not sufficiency.} An interior local extremum of a
differentiable $E$ must satisfy \eqref{app:eq:identity-iff}, so the identity is a
necessary condition for an interior optimum and it identifies the candidates.
It is not sufficient: a point where the elasticities agree may be a local
maximum, and the identity says nothing about a boundary optimum. The smallest
counterexample uses a convex accuracy curve, $a=b=1$ and
$A(k)=0.1+0.05k^{2}$. Then $\Delta(k)=0$ at $k=\sqrt{3}-1\approx0.732$, where
$E\approx13.66$ J/correct, while $E(0)=E(2)=10$ J/correct; the stationary point
is a local maximum, and the hypothesis $A''\le0$ of Proposition~\ref{app:prop:unimodal}
fails. A zero gap is therefore evidence about the local slope of $E$ at a depth,
never by itself evidence that the depth is optimal.

\paragraph{Reading the sign.} Because the prefactor in \eqref{app:eq:proof-second}
is positive, $\operatorname{sign}E'(k)=\operatorname{sign}\Delta(k)$ for $k>0$. A
positive gap means the marginal joule buys less than the marginal correct answer,
so a small increase in depth raises energy per correct answer; a negative gap
means the opposite. This is the interpretation used for the reported gap, and it
is invariant to the depth-common factor $f$ of (A6), because both elasticities
are logarithmic derivatives.

\subsection{Proposition 1: unimodality under concave accuracy}\label{app:prop:unimodal}

\paragraph{Proposition 1.} Assume (A1), (A2), and $A''(k)\le0$ on $(0,K)$; write
$J(k)=a+bk$ and define the stationarity function
\begin{equation}\label{app:eq:N}
  N(k)=bA(k)-J(k)A'(k)=A(k)^{2}E'(k)
  \qquad [\mathrm{J\,query^{-1}\,passage^{-1}}].
\end{equation}
Then:
\begin{enumerate}
  \item[(i)] $N'(k)=-(a+bk)A''(k)\ge0$, so $N$ is nondecreasing on $[0,K]$, and
  strictly increasing if $A''<0$ on $(0,K)$.
  \item[(ii)] $E'=N/A^{2}$ changes sign at most once on $(0,K)$, and any sign
  change is from negative to positive. Hence $E$ is nonincreasing then
  nondecreasing: it is unimodal (quasiconvex) and has no interior local maximum.
  If $A''<0$, $E$ has at most one interior stationary point, and when it exists it
  is the unique global minimiser of $E$ on $[0,K]$.
  \item[(iii)] If $A''<0$, an interior minimiser exists if and only if
  \begin{equation}\label{app:eq:interior}
    bA(0)<aA'(0)
    \quad\text{and}\quad
    bA(K)>J(K)A'(K);
  \end{equation}
  otherwise the minimiser is $k=0$ (if the first inequality fails) or $k=K$ (if
  the second fails).
\end{enumerate}

\paragraph{Proof.} (i) Differentiate \eqref{app:eq:N} using (A1):
\begin{equation}\label{app:eq:Nprime}
  N'(k)=bA'(k)-\left[bA'(k)+(a+bk)A''(k)\right]=-(a+bk)A''(k),
\end{equation}
and $a+bk>0$ by (A1), so $N'\ge0$; if $A''<0$ then $N'>0$.

(ii) By \eqref{app:eq:N}, $E'=N/A^{2}$ and $A^{2}>0$, so $E'$ has the sign of
$N$. A nondecreasing function crosses zero at most once, and when it does the
crossing is from negative to positive; hence there is a point $\bar k$ (possibly
$0$ or $K$, possibly an interval when $N$ is flat at zero) with $E'\le0$ for
$k<\bar k$ and $E'\ge0$ for $k>\bar k$. That is unimodality, and no interior local
maximum is possible. If $A''<0$ then $N$ is strictly increasing, so its zero set
is empty or a single point $k^{*}$. At such a point,
\begin{equation}\label{app:eq:Esecond}
  E''(k^{*})=\frac{N'(k^{*})A(k^{*})^{2}-2A(k^{*})A'(k^{*})N(k^{*})}{A(k^{*})^{4}}
  =\frac{N'(k^{*})}{A(k^{*})^{2}}>0,
\end{equation}
because $N(k^{*})=0$; the point is a strict local minimum, and by the sign pattern
it is the unique global minimiser on $[0,K]$.

(iii) If $A''<0$ then $N$ is continuous and strictly increasing. An interior zero
exists if and only if $N(0)<0<N(K)$; substituting \eqref{app:eq:N} gives
$bA(0)-aA'(0)<0$ and $bA(K)-(a+bK)A'(K)>0$, which is \eqref{app:eq:interior}.
If $N(0)\ge0$ then $N>0$ on $(0,K]$ and $E$ is nondecreasing, so the minimiser is
$k=0$; if $N(K)\le0$ then $N<0$ on $[0,K)$ and $E$ is nonincreasing, so the
minimiser is $k=K$. \hfill$\square$

\paragraph{Two remarks.} First, the proof never uses $A'\ge0$: monotonicity of
$N$ depends only on the sign of $A''$. The proposition therefore also covers a
concave decline after an accuracy peak, which is the qualitative shape of the two
non-monotone P4 arms (a peak followed by a decline), although concavity of the
decline is not tested there. Second, the hypothesis is testable on the grid: the discrete
analogue of $A''\le0$ is that the secant slopes between adjacent grid depths are
non-increasing, which is what the companion implementation checks. The hypothesis
is not a description of every accuracy curve; where it fails, Proposition~1 gives
no global statement and the identity remains local.

\paragraph{Existence in elasticity form.} Conditions \eqref{app:eq:interior} can
be read as bounds on the elasticities. The first is
$b/a<A'(0)/A(0)$, i.e.\ the relative accuracy gain per passage at zero must exceed
the relative energy cost per passage. The second is
$\varepsilon_{A}(K)<\varepsilon_{J}(K)=bK/(a+bK)<1$, i.e.\ the accuracy
elasticity at the top of the grid must be below the energy elasticity. Both are
properties of the fitted curves, not of the decision rule, and both are evaluated
without the accuracy gate of (A3).

\subsection{Corollary 1: comparative statics of the optimum}\label{app:cor:statics}

\paragraph{Corollary 1.} Assume (A1), (A2), $A'(k^{*})>0$ and $A''(k^{*})<0$ at an
interior stationary point $k^{*}=k^{*}(a,b)$ of $E$. Then $k^{*}$ is continuously
differentiable in $(a,b)$ in a neighbourhood, and
\begin{equation}\label{app:eq:statics}
  \frac{\partial k^{*}}{\partial a}
  =\frac{-A'(k^{*})}{(a+bk^{*})A''(k^{*})}>0,
  \qquad
  \frac{\partial k^{*}}{\partial b}
  =\frac{A(k^{*})-k^{*}A'(k^{*})}{(a+bk^{*})A''(k^{*})}
  =-\frac{aA(k^{*})}{(a+bk^{*})^{2}\left(-A''(k^{*})\right)}<0 .
\end{equation}

\paragraph{Proof.} By Lemma~1 the stationary point solves
$F(k,a,b)=bA(k)-(a+bk)A'(k)=0$, and $F$ is continuously differentiable under
(A2). Its partials are
\begin{equation}\label{app:eq:Fpartials}
  F_{k}=-(a+bk)A''(k),
  \qquad
  F_{a}=-A'(k),
  \qquad
  F_{b}=A(k)-kA'(k).
\end{equation}
At $k^{*}$, $F_{k}=-(a+bk^{*})A''(k^{*})>0$ because $A''(k^{*})<0$; this is also
the second-order condition \eqref{app:eq:Esecond}. The implicit function theorem
therefore applies, and
\begin{equation}\label{app:eq:IFT}
  \frac{\partial k^{*}}{\partial a}=-\frac{F_{a}}{F_{k}},
  \qquad
  \frac{\partial k^{*}}{\partial b}=-\frac{F_{b}}{F_{k}} .
\end{equation}
Substituting \eqref{app:eq:Fpartials} gives the first equality in each line of
\eqref{app:eq:statics}. The first sign is immediate since $A'(k^{*})>0$ and the
denominator is negative. For the second, apply the identity
\eqref{app:eq:identity-iff} at $k^{*}$:
$k^{*}A'(k^{*})/A(k^{*})=\varepsilon_{J}(k^{*})=bk^{*}/(a+bk^{*})<1$, hence
\begin{equation}\label{app:eq:Fbpositive}
  A(k^{*})-k^{*}A'(k^{*})
  =A(k^{*})\left(1-\varepsilon_{J}(k^{*})\right)
  =\frac{aA(k^{*})}{a+bk^{*}}>0,
\end{equation}
and the denominator of the second line is negative, so
$\partial k^{*}/\partial b<0$. \hfill$\square$

\paragraph{Reading.} A larger fixed cost per query makes depth relatively cheap
and pushes the unconstrained optimum deeper; a larger per-passage cost makes
depth relatively expensive and pushes it shallower. The corollary is local and
applies only when the accuracy gate is slack: when the gate binds, the policy
depth is the smallest feasible grid depth and is locally insensitive to $(a,b)$,
so the derivative is zero in the one-sided sense at that corner. It is also
silent about the magnitude, because $A'(k^{*})$ and $A''(k^{*})$ are estimated
from the same finite sample as $A(k^{*})$; the sign is a prediction, the size is
not.

\subsection{Corollary 2: when the optimum is the boundary}\label{app:cor:boundary}

\paragraph{Corollary 2.} Assume (A1) and the stylized coverage family
\begin{equation}\label{app:eq:coverage}
  A(k)=1-(1-a_{0})(1-p)^{k},
  \qquad
  a_{0}=A(0)\in(0,1),
  \qquad
  p\in(0,1),
\end{equation}
so that $A(k)\in(a_{0},1)$ and (A2) holds on $[0,\infty)$. Let
$L=-\ln(1-p)>0$. Then $A'(k)=L\left(1-A(k)\right)>0$ and
$A''(k)=-L^{2}\left(1-A(k)\right)<0$, Proposition~\ref{app:prop:unimodal}
applies, and
\begin{equation}\label{app:eq:boundary}
  k^{*}=0
  \iff
  \frac{b}{a}\ge\frac{(1-a_{0})\left(-\ln(1-p)\right)}{a_{0}}
  =\frac{1-a_{0}}{a_{0}}\,L .
\end{equation}
When the inequality is reversed, the unique interior optimum solves
\begin{equation}\label{app:eq:coverage-optimum}
  (1-a_{0})(1-p)^{k^{*}}=\frac{b}{aL+b+bLk^{*}} .
\end{equation}

\paragraph{Proof.} Differentiating \eqref{app:eq:coverage} gives
$A'(k)=-(1-a_{0})(1-p)^{k}\ln(1-p)=L\left(1-A(k)\right)$ and
$A''(k)=-L^{2}\left(1-A(k)\right)<0$, so $A$ is strictly increasing and strictly
concave and (A2) holds. Evaluate \eqref{app:eq:N} at zero:
\begin{equation}\label{app:eq:N0}
  N(0)=bA(0)-aA'(0)=ba_{0}-a(1-a_{0})L
  =a_{0}\left[b-a\frac{1-a_{0}}{a_{0}}L\right].
\end{equation}
By Proposition~\ref{app:prop:unimodal}(i), $N$ is strictly increasing. As
$k\to\infty$, $A(k)\to1$ and $(a+bk)(1-A(k))\to0$ because the exponential decay
dominates the linear factor, so
\begin{equation}\label{app:eq:Nlimit}
  N(k)=bA(k)-(a+bk)L\left(1-A(k)\right)\longrightarrow b>0 .
\end{equation}
If $N(0)\ge0$, then $N>0$ on $(0,\infty)$ and $E$ is strictly increasing, so
$k^{*}=0$; by \eqref{app:eq:N0} this is exactly the inequality in
\eqref{app:eq:boundary}. If $N(0)<0$, then \eqref{app:eq:Nlimit} and strict
monotonicity give a unique zero $k^{*}\in(0,\infty)$, which is the unique
minimiser by Proposition~\ref{app:prop:unimodal}(ii). Setting $N(k^{*})=0$ and
writing $s=1-A(k^{*})=(1-a_{0})(1-p)^{k^{*}}$ gives
$b(1-s)=(a+bk^{*})Ls$, i.e.\ $b=s\left(aL+b+bLk^{*}\right)$, which is
\eqref{app:eq:coverage-optimum}. \hfill$\square$

\paragraph{Scope.} The corollary is a statement about one stylized family, not
about the observed curves. The P4 validation rejects the family as a general
description of the four arms ($R^{2}$ between $0.097$ and $0.879$), so its role
is to show that a boundary optimum is a possible outcome of diminishing returns,
and to give the threshold in closed form. It does not certify any observed
$k^{*}$, and it is not used to choose a depth.

\subsection{Proposition 2: the price of static depth}\label{app:prop:price}

\paragraph{Proposition 2.} Assume (A1). Let $\pi$ be any per-question depth
assignment with values in $[0,K]$, with population energy
$J_{\pi}=\mathbb{E}_{q}[J(\pi(q))]$ and accuracy
$A_{\pi}=\mathbb{E}_{q}[A(\pi(q))]$, and let $k_{s}$ be a static depth with
accuracy $A_{s}=A(k_{s})$. If $A_{\pi}\ge A_{s}$, then
\begin{equation}\label{app:eq:ceiling}
  \frac{E_{\pi}}{E_{s}}
  =\frac{J_{\pi}}{J(k_{s})}\cdot\frac{A_{s}}{A_{\pi}}
  \le\frac{a+bK}{a}
  =\frac{1}{1-\varepsilon_{J}(K)} ,
\end{equation}
and the largest saving any such assignment can deliver is the variable energy
share at the reference depth,
\begin{equation}\label{app:eq:ceiling-saving}
  1-\frac{E_{\pi}}{E_{s}}\le\frac{bK}{a+bK}=\varepsilon_{J}(K)<1 .
\end{equation}
The bound is sharp. For $\varepsilon\in(0,1)$ let a fraction $1-\varepsilon$ of
questions be easy, answered correctly at every depth including $k=0$, and a
fraction $\varepsilon$ be hard, answered correctly only at depths $k\ge K$. Then
the best static depth at accuracy $1$ is $k_{s}=K$, the hindsight oracle
$\pi_{\mathrm{or}}$ assigns easy questions to $0$ and hard questions to $K$, and
\begin{equation}\label{app:eq:rho}
  \rho(\varepsilon)
  =\frac{E_{\mathrm{static}}}{E_{\mathrm{oracle}}}
  =\frac{a+bK}{a+\varepsilon bK}
  \;\longrightarrow\;
  1+\frac{bK}{a}
  =\frac{1}{1-\varepsilon_{J}(K)}
  \qquad(\varepsilon\to0).
\end{equation}
The ratio is written here as static over oracle, so that the quantity is the
\emph{price} of fixing one depth; written in the opposite order (oracle over
static) the same limit is its reciprocal $a/(a+bK)<1$.

\paragraph{Proof.} By (A1), $J(\pi(q))\in[a,a+bK]$ for every question, so
$a\le J_{\pi}\le a+bK$ and $J(k_{s})\ge a$; hence
$J_{\pi}/J(k_{s})\le(a+bK)/a$. Since $A_{\pi}\ge A_{s}>0$, the factor
$A_{s}/A_{\pi}$ is at most one, which proves \eqref{app:eq:ceiling}. For
\eqref{app:eq:ceiling-saving}, subtract the reciprocal:
$1-a/(a+bK)=bK/(a+bK)=\varepsilon_{J}(K)$ by (A1).

For sharpness, let $q$ be easy with probability $1-\varepsilon$ and hard with
probability $\varepsilon$. Every static depth $k<K$ answers only the easy
questions, so its accuracy is $1-\varepsilon<1$; the only static depths with
accuracy $1$ are $k\ge K$, and among them $J$ is minimised at $k=K$. Hence
$E_{\mathrm{static}}=a+bK$. The oracle answers every question correctly, so its
accuracy is $1$ and its energy is
$(1-\varepsilon)a+\varepsilon(a+bK)=a+\varepsilon bK$, giving
$E_{\mathrm{oracle}}=a+\varepsilon bK$ and the first equality in
\eqref{app:eq:rho}. The limit is immediate, and the last equality repeats
$bK/(a+bK)=\varepsilon_{J}(K)$. \hfill$\square$

\paragraph{Scope.} Proposition~\ref{app:prop:price} is an upper bound with a
sharpness example, not a lower bound on the value of adaptivity and not a claim
about the observed workload. Two consequences matter for the paper. First, the
ceiling is a property of the affine energy model alone: it caps \emph{any}
allocation, learned or oracle, at matched accuracy, and it is the formal reason
the oracle gain is reported as a fraction of the static energy rather than as an
unbounded promise. Second, the example shows that the ceiling is approached when
a small hard stratum forces the whole workload to the deepest setting; when the
stratum is large ($\varepsilon\to1$), the ratio tends to one and adaptivity is
worth nothing. The empirical matched-oracle gain reported for P3 is far below the
P1 ceiling of $\varepsilon_{J}(32)=0.7517$, as it must be, and the pilot oracle
figure of $78.8\%$ is below the pilot ceiling
$bK/(a+bK)=0.8227$ computed from $a=15.74$ and $b=2.283$; both are consistency
checks on registered quantities, not new evidence.

\subsection{Proposition 3: the cost of certification}\label{app:prop:cert}

Certification has two separate costs: the energy of the experiment (part a) and
the number of questions needed for the statistical statement (part b). They are
stated separately because the first is a deterministic break-even and the second
is a sampling requirement, and only their product is the cost of a certified
depth.

\paragraph{Proposition 3(a): economic break-even.} Assume (A1), a deployment of
$N$ queries drawn from the same distribution as the experiment, one depth used
for all $N$ queries, no discounting, and a one-off certification cost
$C_{\mathrm{cert}}>0$ in joules. Let $k_{0}=K$ be the reference depth and
$k^{*}$ the candidate. Define the per-query break-even and the
per-correct-answer break-even by
\begin{equation}\label{app:eq:nstar}
  N^{*}_{J}=\frac{C_{\mathrm{cert}}}{J(k_{0})-J(k^{*})},
  \qquad
  N^{*}_{A}=\frac{C_{\mathrm{cert}}}{A(k^{*})\left[E(k_{0})-E(k^{*})\right]},
\end{equation}
both in queries and both defined only when the denominator is positive. Then
\begin{equation}\label{app:eq:nstar-identity}
  A(k^{*})\left[E(k_{0})-E(k^{*})\right]
  =\frac{A(k^{*})}{A(k_{0})}J(k_{0})-J(k^{*}),
\end{equation}
so $N^{*}_{A}=N^{*}_{J}$ if and only if $A(k^{*})=A(k_{0})$, and otherwise
\begin{equation}\label{app:eq:beta}
  \beta:=\frac{N^{*}_{A}}{N^{*}_{J}}
  =\frac{J(k_{0})-J(k^{*})}{\dfrac{A(k^{*})}{A(k_{0})}J(k_{0})-J(k^{*})},
\end{equation}
with $\beta\ge1$ if and only if $A(k^{*})\le A(k_{0})$.

\paragraph{Proof.} Equation \eqref{app:eq:nstar-identity} is
$A(k^{*})E(k_{0})-A(k^{*})E(k^{*})=A(k^{*})J(k_{0})/A(k_{0})-J(k^{*})$, using
$A(k^{*})E(k^{*})=J(k^{*})$. Dividing the two lines of \eqref{app:eq:nstar}
gives \eqref{app:eq:beta}. For the sign, compare denominators:
\begin{equation}\label{app:eq:beta-sign}
  \left[J(k_{0})-J(k^{*})\right]
  -\left[\frac{A(k^{*})}{A(k_{0})}J(k_{0})-J(k^{*})\right]
  =J(k_{0})\left(1-\frac{A(k^{*})}{A(k_{0})}\right),
\end{equation}
which is non-negative exactly when $A(k^{*})\le A(k_{0})$. Both denominators are
positive by hypothesis, so the inequality passes to $\beta$.
\hfill$\square$

\paragraph{Interpretation and the approximation.} $N^{*}_{J}$ is the number of
deployment queries at which the energy saved by switching depth equals the
certificate: after $N$ queries at $k^{*}$ instead of $k_{0}$ the saving is
$N\left[J(k_{0})-J(k^{*})\right]$, and it repays $C_{\mathrm{cert}}$ at
$N=N^{*}_{J}$. $N^{*}_{A}$ answers a different question: it values the candidate's
correct-answer output at the per-correct-answer saving $E(k_{0})-E(k^{*})$, so $N$
candidate-depth queries produce $N A(k^{*})$ correct answers and a cumulative
saving $N A(k^{*})\left[E(k_{0})-E(k^{*})\right]$, which equals
$C_{\mathrm{cert}}$ at $N=N^{*}_{A}$. The two coincide only at
equal accuracy. When $A(k^{*})<A(k_{0})$, the paper's per-correct-answer formula
is therefore an approximation to the deployment-queries reading, and it errs on
the conservative side ($\beta>1$): it reports a longer break-even because it also
requires the candidate to reproduce the correct-answer output that its lower
accuracy would otherwise lose. The direction is known but the formula is not
exact, which is why both numbers are reported rather than one.

\paragraph{P1 diagnostic.} With the registered P1 values
$C_{\mathrm{cert}}=311{,}992.064$ J, $J(4)=36.2537$ J/query,
$J(32)=94.9450$ J/query, $A(4)=0.450$, and $A(32)=0.493$, one obtains
$N^{*}_{J}=5{,}315.8$ queries, $N^{*}_{A}=6{,}189.1$ queries, and
$\beta=1.1643$. The certification cost is itself a sum over the grid,
$C_{\mathrm{cert}}=n\sum_{k\in\mathcal{K}}J(k)=1{,}000\times311.992$ J, so the
sweep energy and the reported break-even use the same accounting.

\paragraph{Proposition 3(b): statistical sample complexity.} Assume (A5), let
$m$ be the accuracy margin (written $\delta$ in the main text, and renamed here
to avoid a clash with the error probability), and let
$\Delta=A(k_{0})-A(k^{*})>0$. For the paired differences
$D_{q}=C_{qk_{0}}-C_{qk^{*}}\in\{-1,0,1\}$ with
$\mathbb{E}[D_{q}]=\Delta$, put
$t_{n}=\sqrt{2\ln(1/\alpha)/n}$ and declare non-inferiority of $k^{*}$ when
\begin{equation}\label{app:eq:declare}
  \bar D+t_{n}<m .
\end{equation}
Then:
\begin{enumerate}
  \item[(i)] \eqref{app:eq:declare} is a level-$\alpha$ test of
  $H_{0}:\Delta\ge m$.
  \item[(ii)] At a true difference $\Delta=m_{1}<m$ the probability of failing to
  declare non-inferiority is at most
  \begin{equation}\label{app:eq:hoeffding}
    \Pr\left(\text{fail to declare}\mid \Delta=m_{1}\right)
    \le\exp\left(-\frac{n\left(\left(m-m_{1}\right)-t_{n}\right)^{2}}{2}\right),
  \end{equation}
  so $n\ge8\max\left(\ln(1/\alpha),\ln(1/\delta_{\mathrm{err}})\right)/(m-m_{1})^{2}$
  paired questions suffice for a failure probability at most
  $\delta_{\mathrm{err}}$.
  \item[(iii)] Selecting the more accurate of two depths at a true gap $\Delta$
  needs $n\ge2\ln(1/\delta_{\mathrm{err}})/\Delta^{2}$, and selecting among $d$
  depths whose smallest gap to the optimum is $\Delta_{\min}$ needs
  $n\ge2\ln\left((d-1)/\delta_{\mathrm{err}}\right)/\Delta_{\min}^{2}$ by a union
  bound; in both cases the correct depth is selected with probability at least
  $1-\delta_{\mathrm{err}}$.
\end{enumerate}

\paragraph{Proof sketch.} $D_{q}$ is bounded in $[-1,1]$ and has range $2$, so
Hoeffding's inequality gives, for every $u>0$,
\begin{equation}\label{app:eq:two-tails}
  \Pr\left(\bar D-\Delta\le-u\right)\le\exp\left(-\frac{nu^{2}}{2}\right),
  \qquad
  \Pr\left(\bar D-\Delta\ge u\right)\le\exp\left(-\frac{nu^{2}}{2}\right).
\end{equation}
(i) Under $H_{0}$, $\Delta\ge m$, so declaring non-inferiority implies
$\bar D-\Delta<-t_{n}$, which has probability at most
$\exp(-nt_{n}^{2}/2)=\alpha$.

(ii) Failure to declare means $\bar D+t_{n}\ge m$, hence
$\bar D-\Delta\ge(m-m_{1})-t_{n}$, and the upper tail of
\eqref{app:eq:two-tails} applies. If $t_{n}\le(m-m_{1})/2$, which holds once
$n\ge8\ln(1/\alpha)/(m-m_{1})^{2}$, the exponent is at most
$-n(m-m_{1})^{2}/8$, and the stated $n$ makes it at most
$\ln(1/\delta_{\mathrm{err}})$.

(iii) The selection error is
$\Pr(\bar D\le0)=\Pr(\bar D-\Delta\le-\Delta)\le\exp(-n\Delta^{2}/2)$, and the
union bound over the $d-1$ alternatives gives the second statement.
\hfill$\square$

\paragraph{Scope and consequences.} The bound is distribution-free and makes no
normal approximation, but it is deliberately conservative: it ignores the
pairing, i.e.\ the discordance rate, which is what a McNemar calculation uses.
The tighter normal-approximation requirement scales like
$\psi/(m-m_{1})^{2}$ at discordance $\psi$, so at a fixed difference the observed
P1 discordance of $13.1\%$ requires roughly $1.6$ times the pairs that the
planned $8\%$ implied, and the registered frontier puts the 5 pp requirement at
$n=7{,}138$ pairs. Three consequences are used in the main text. First, the
sample complexity diverges as $m_{1}\uparrow m$: a margin that the data show is
several percentage points wide cannot be resolved by a design powered at a
near-zero difference, which is the formal content of the P1 precision failure.
Second, the energy side does not drive the requirement: per-query energy has a P1
repeat CV of $2.81\%$ at $k=8$, so the relative standard error of a mean energy
at $n=1{,}000$ is about $2.81\%/\sqrt{1{,}000}<0.1\%$, while the accuracy
contrast at 5 pp remains unresolved at the same $n$. The accuracy side is the
binding one. Third, the certification cost is linear in the statistical
requirement, because a sweep of $n$ questions costs
$C_{\mathrm{cert}}(n)=n\,c_{\mathrm{sweep}}$ with
$c_{\mathrm{sweep}}=\sum_{k\in\mathcal{K}}J(k)$ in J/query, so
\begin{equation}\label{app:eq:nstar-n}
  N^{*}_{J}(n)=\frac{n\,c_{\mathrm{sweep}}}{J(k_{0})-J(k^{*})},
\end{equation}
which grows linearly in $n$. The statistical requirement and the economic
break-even therefore cannot be chosen independently, and a depth that needs a
large experiment needs a proportionally larger deployment to justify it.

\subsection{Grid resolution and the H3 diagnostic}\label{app:grid}

The identity of Lemma~1 is a statement about a derivative. On the grid two
approximations intervene, and separating them is what makes the H3 result
readable.

\emph{First, the derivative is a secant.} The reported accuracy elasticity at a
grid depth is a finite difference over the neighbouring depths, so the identity
at a grid point is the statement that the secant elasticity equals the affine
elasticity. At $k=4$ the central difference uses $k=2$ and $k=8$, a step of three
passages; at $k=8$ it uses $k=4$ and $k=16$, a step of six; at the edges the
difference is one-sided, with step one at $k=1$ and step sixteen at $k=32$. A
central difference approximates $A'(k)$ to second order in the step and a
one-sided difference to first order, so a grid gap is a discretisation of
$\Delta(k)$, not a different quantity.

\emph{Second, the decision object is constrained.} The policy depth of (A3) is the
energy minimiser inside the feasible set $\mathcal{F}_{m}$, whereas $k^{*}$ in
Lemma~1 is the unconstrained stationary point of $E$. The two objects coincide
only when the accuracy gate is slack. In P1 the gate binds: the 5 pp rule selects
$k=4$ because it is the smallest feasible depth, not because $E$ is stationary
there. Testing the identity at the gate-selected depth is therefore a test of
local stationarity at a point that the theory does not expect to be stationary,
and the main text reports it as such.

\emph{The P1 numbers.} With $a=24.0759$ J/query and
$b=2.2772$ J per passage per query, the affine elasticity is
$\varepsilon_{J}(4)=0.2745$, the secant elasticity is
$\varepsilon_{A}(4)=0.0741$, and the gap is $\Delta(4)=+0.2004$ with a 95\%
interval of $[0.1677,0.2320]$ that excludes zero. The positive sign says that
$E$ is increasing at $k=4$: the last passages bought at that depth cost more
energy per correct answer than they return. The monotone-cubic interpolant of the
same seven means has its unconstrained stationary point at $k^{*}=1.2144$ with
$E=67.1408$ J/correct and a gap of $-1.9\times10^{-7}$, which is zero to numerical
precision. Since $1.2144<4$, a positive gap at $k=4$ is exactly what the theory
requires. The H3 test asked whether the gate-selected depth is stationary, and
the answer is no; that is a statement about the selected depth, not a failure of
the identity.

\emph{What the interpolation does and does not establish.} The near-zero gap at
$k^{*}=1.2144$ is close to a consistency check on the solver: any smooth
interpolant, minimised numerically, will report $E'\approx0$ at its own
minimiser, so the zero is not independent evidence for the location. The
independent content is the sign at $k=4$, which is estimated from the data and
has an interval that excludes zero, and the flatness of $E$ between $k=1$ and
$k=2$: the two values are $73.71$ and $73.33$ J/correct, a difference of $0.5\%$,
with overlapping percentile intervals $[68.2,79.9]$ and $[68.1,79.3]$. The
existence of an interior unconstrained optimum is therefore supported on this
arm, while its location inside $[1,2]$ is not identified. The defensible
statement is that $k=4$ is over-retrieval relative to an unconstrained optimum
near $k=1.2$, not that $1.214$ is the optimal depth. Depth $0$ is excluded from
the identity because $\varepsilon$ is degenerate there and because $A(0)=0.001$
rests on a single correct answer.

\subsection{What is not claimed}\label{app:notclaimed}

The appendix is deliberately narrower than a general theory of retrieval depth.
It does not claim the following.

\begin{itemize}
  \item \emph{A universal $k^{*}$.} Every number quoted here is a property of one
  reader, one retriever, one corpus, one grid, and one hardware condition. The
  comparative statics are signs, not magnitudes, and they are derived for the
  unconstrained ratio.
  \item \emph{That accuracy is concave in general.} Proposition~\ref{app:prop:unimodal}
  is conditional. The P4 validation rejects the coverage family on three of four
  arms and finds non-monotone curves on two, so the hypothesis is tested where it
  is used and is not a maintained assumption about RAG.
  \item \emph{That a zero gap certifies optimality.} The identity is necessary,
  not sufficient, as the convex counterexample in Section~\ref{app:lemma1} shows.
  A gap interval that contains zero is consistent with stationarity; it does not
  establish a minimum, and it cannot speak for a boundary depth.
  \item \emph{That the grid argmin is the continuous argmin.} The grid is coarse
  and uneven, the interpolant is a modelling choice, and the policy object is the
  constrained minimiser of (A3). The two optima are reported separately and are
  never substituted for one another.
  \item \emph{That the ceiling of Proposition~\ref{app:prop:price} is attained by
  a deployable policy.} It is a bound with a sharpness example in a stylized
  mixture. A learned policy is a separate object that must be trained, selected,
  and evaluated on disjoint splits, and the P3 result shows that a realizable
  policy can capture none of the gap.
  \item \emph{That $N^{*}$ is a deployment guarantee.} It is a deterministic
  break-even: no discounting, no uncertainty in the saving, no cost of capital,
  and a certificate that is assumed to settle the question once paid.
  \item \emph{Any claim outside the measured condition.} No carbon claim, no
  claim about other hardware, readers, or retrievers, and no claim that the
  exploratory arms enter the theory.
  \item \emph{Any re-adjudication.} The appendix proves statements; it does not
  change the status of H3 or of any gate. H3 remains adjudicated at $k=4$ and
  remains unsupported.
\end{itemize}
